\documentclass[11pt,letterpaper]{article}
\usepackage[margin=0.82in]{geometry}
\usepackage[T1]{fontenc}
\usepackage{lmodern,amsmath,amssymb,mathtools,booktabs,longtable,array,graphicx,microtype,xcolor,enumitem,fancyhdr,hyperref}
\usepackage[most]{tcolorbox}

\hypersetup{
    colorlinks=true,
    linkcolor=blue!55!black,
    urlcolor=blue!55!black,
    pdftitle={MAT235 Week 2 Answers and Practice}
}

\setlength{\headheight}{14pt}
\pagestyle{fancy}
\fancyhf{}
\lhead{MAT235 --- Week 2 Answers \& Practice}
\rhead{Worked Solutions, Textbook, \& Self-Test}
\cfoot{\thepage}

\setcounter{tocdepth}{2}
\setlength{\parskip}{4pt}
\setlength{\parindent}{0pt}

\newcommand{\fig}[2]{\begin{center}\includegraphics[width=#1\linewidth]{graphs/#2}\end{center}}
\newcommand{\R}{\mathbb{R}}

\begin{document}
\hypersetup{pageanchor=false}

\begin{titlepage}
\centering
\vspace*{1.2in}
{\Huge\bfseries MAT235 Week 2\\[12pt] Answers and Practice\par}
\vspace{18pt}
{\LARGE Complete Solutions to Problem Types, Weekly Source Questions, Suggested Textbook Exercises, and Original Self-Test\par}
\vspace{22pt}
{\large Comprehensive Reference Companion to \texttt{Study-Notes.tex}\par}
\vfill
{\large Sections 12.2 and 12.3 with Supplementary Review and Previews\par}
\vspace{10pt}
{\small Verified Against 8th-Edition Textbook, Lecture Slides, and Handwritten Course Materials\par}
\end{titlepage}

\hypersetup{pageanchor=true}
\pagenumbering{roman}
\tableofcontents
\clearpage
\pagenumbering{arabic}

% =========================================================================
% SECTION A: PROBLEM TYPE SOLUTIONS
% =========================================================================
\section{Problem Type Solutions}\label{sec:pt-solutions}
This section provides complete worked solutions and final answers for every Problem Type teaching example introduced in \texttt{Study-Notes.tex}.

\subsection{Problem Type PT-01: Recognize and Describe a Surface from an Equation}
\textbf{Question reference.} See PT-01 in \texttt{Study-Notes.tex}.

\textbf{Worked solution.}
\begin{enumerate}[label=(\alph*),leftmargin=*]
    \item For $z = 1 + x - y$:
    Rearranging into standard linear form gives $x - y - z = -1$.
    \begin{itemize}[nosep]
        \item \emph{Landmarks:} Setting $y=0, z=0$ gives $x$-intercept $(-1,0,0)$; setting $x=0, z=0$ gives $y$-intercept $(0,1,0)$; setting $x=0, y=0$ gives $z$-intercept $(0,0,1)$.
        \item \emph{Traces:} In the plane $x=0$, the trace is the line $z = 1 - y$; in $y=0$, the line $z = 1 + x$; in $z=0$, the line $y = 1 + x$.
        \item \emph{Function test:} For every $(x,y) \in \R^2$, $z$ is uniquely given by $1+x-y$. Thus, it satisfies the vertical-line test everywhere.
        \item \emph{Orientation and extent:} The surface is an infinite, unbounded plane that slopes upward in the direction of increasing $x$ and downward in the direction of increasing $y$.
    \end{itemize}
    \item For $x^2 + y^2 + z^2 = 4$:
    This matches the standard sphere equation $(x-x_0)^2 + (y-y_0)^2 + (z-z_0)^2 = r^2$ with center $(0,0,0)$ and radius $r = 2$.
    \begin{itemize}[nosep]
        \item \emph{Landmarks:} Axis intercepts are $(\pm 2, 0, 0)$, $(0, \pm 2, 0)$, and $(0, 0, \pm 2)$.
        \item \emph{Traces:} Setting $x=0$ gives the circle $y^2+z^2=4$ in the $yz$-plane; setting $y=0$ gives $x^2+z^2=4$ in the $xz$-plane; setting $z=0$ gives $x^2+y^2=4$ in the $xy$-plane.
        \item \emph{Function test:} At $(x,y)=(0,0)$, $z^2 = 4 \implies z = \pm 2$. The vertical line through the origin intersects the surface at $(0,0,2)$ and $(0,0,-2)$. Thus, the full sphere fails the vertical-line test and is not the graph of a single function.
        \item \emph{Orientation and extent:} The surface is completely bounded within $[-2,2] \times [-2,2] \times [-2,2]$. It is the union of two function graphs $z = \pm\sqrt{4-x^2-y^2}$ over the closed disk $x^2+y^2 \le 4$.
    \end{itemize}
\end{enumerate}

\textbf{Answer.}
(a) The equation $z = 1 + x - y$ represents an infinite plane with intercepts $(-1,0,0)$, $(0,1,0)$, and $(0,0,1)$; it is the graph of a single-valued function $z = f(x,y)$ on all of $\R^2$.
(b) The equation $x^2 + y^2 + z^2 = 4$ represents a bounded sphere of radius $2$ centered at the origin; it fails the vertical-line test and cannot be written as a single function $z = f(x,y)$, but decomposes into upper and lower hemispheres over the disk $x^2+y^2 \le 4$.

\subsection{Problem Type PT-02: Test a Surface as a Function $z=f(x,y)$ and Determine Exact Domain}
\textbf{Question reference.} See PT-02 in \texttt{Study-Notes.tex}.

\textbf{Worked solution.}
\begin{enumerate}[label=(\alph*),leftmargin=*]
    \item Evaluate the equation $x^2 + y^2 + z^2 = 1$ at $(x,y) = (0,0)$:
    \[ 0^2 + 0^2 + z^2 = 1 \implies z^2 = 1 \implies z = 1 \quad \text{or} \quad z = -1. \]
    The vertical line $x=0, y=0$ intersects the sphere at two distinct spatial points, $(0,0,1)$ and $(0,0,-1)$. By definition, a single-valued function assigns exactly one output to each input pair. Hence, the whole sphere fails the vertical-line test.
    \item Solving $x^2 + y^2 + z^2 = 1$ for $z$ yields:
    \[ z^2 = 1 - x^2 - y^2 \implies z = \pm\sqrt{1 - x^2 - y^2}. \]
    This gives two single-valued functions:
    \[ f_1(x,y) = +\sqrt{1 - x^2 - y^2} \quad \text{(upper hemisphere, } z \ge 0\text{)}, \]
    \[ f_2(x,y) = -\sqrt{1 - x^2 - y^2} \quad \text{(lower hemisphere, } z \le 0\text{)}. \]
    \item For the square root to produce real values, the radicand must be non-negative:
    \[ 1 - x^2 - y^2 \ge 0 \iff x^2 + y^2 \le 1. \]
    In set-builder notation, the exact domain is $D = \{(x,y) \in \R^2 : x^2 + y^2 \le 1\}$, which is a closed circular disk of radius $1$. Specifying $[-1,1] \times [-1,1]$ is incorrect because points in the corners of the square (such as $(1,1)$) have $x^2 + y^2 = 2 > 1$, resulting in a strictly negative radicand $1 - 2 = -1$ where the square root is not real.
\end{enumerate}

\textbf{Answer.}
(a) The vertical line $x=0, y=0$ intersects the sphere at $(0,0,1)$ and $(0,0,-1)$, failing the vertical-line test.
(b) Upper hemisphere $z = \sqrt{1 - x^2 - y^2}$; lower hemisphere $z = -\sqrt{1 - x^2 - y^2}$.
(c) The exact domain is the closed unit disk $D = \{(x,y) \in \R^2 : x^2 + y^2 \le 1\}$.

\subsection{Problem Type PT-03: Transform a Known Surface Graph via Shifts, Reflections, and Scaling}
\textbf{Question reference.} See PT-03 in \texttt{Study-Notes.tex}.

\textbf{Worked solution.}
\begin{enumerate}[label=(\alph*),leftmargin=*]
    \item $z = (x-1)^2 + (y-2)^2$:
    Replacing $x \mapsto x-1$ and $y \mapsto y-2$ translates the parent surface $z = x^2+y^2$ horizontally by vector $(+1, +2, 0)$. The base vertex at the origin shifts to $(1,2,0)$. The surface is an upward circular paraboloid with vertical axis of symmetry given by the line $x=1, y=2$. Since squared terms are non-negative, the range is $[0, \infty)$.
    \item $z = 5 - x^2 - y^2 = 5 - (x^2+y^2)$:
    The parent paraboloid is first reflected vertically across the $xy$-plane ($z \mapsto -z$, giving $z = -x^2-y^2$) and then translated upward by $5$ units along the $z$-axis ($z \mapsto z+5$). The surface is a downward-opening circular paraboloid. The vertex is $(0,0,5)$, the axis of symmetry is the $z$-axis ($x=0, y=0$), and the range is $(-\infty, 5]$.
    \item $z = e^{-(x^2+y^2)}$:
    Let $r = \sqrt{x^2+y^2}$. Then $z = e^{-r^2}$. Because $z$ depends solely on $r$, the graph has circular rotational symmetry around the $z$-axis. Since $r^2 \ge 0$, the maximum value is $e^0 = 1$, attained at $(0,0,1)$. For all $(x,y)$, $e^{-r^2} > 0$. As distance from the origin $r \to \infty$, $z = e^{-r^2} \to 0$. The graph is an infinite, circularly symmetric Gaussian bell surface with peak $(0,0,1)$, asymptotic to the $xy$-plane ($z=0$), with range $(0,1]$.
\end{enumerate}

\textbf{Answer.}
(a) Upward circular paraboloid translated by $(1,2,0)$, vertex $(1,2,0)$, axis $x=1, y=2$, range $[0, \infty)$.
(b) Downward circular paraboloid reflected across $xy$-plane and raised by $5$, vertex $(0,0,5)$, axis $x=0, y=0$, range $(-\infty, 5]$.
(c) Rotationally symmetric Gaussian bell surface, peak $(0,0,1)$, horizontal asymptote $z=0$, range $(0,1]$.

\subsection{Problem Type PT-04: Determine and Sketch Parabolic and Hyperbolic Vertical Traces}
\textbf{Question reference.} See PT-04 in \texttt{Study-Notes.tex}.

\textbf{Worked solution.}
\begin{enumerate}[label=(\alph*),leftmargin=*]
    \item In the vertical planes $y = b$, we substitute $y = b$ into $z = x^2 - y^2$:
    \[ z = x^2 - b^2. \]
    This describes an upward-opening parabola in the plane $y = b$. In this plane, coordinates are $x$ (horizontal) and $z$ (vertical). The vertex occurs at $x = 0$, giving coordinates $(x,y,z) = (0, b, -b^2)$.
    \begin{itemize}[nosep]
        \item For $y = -1$: trace is $z = x^2 - 1$ in plane $y = -1$, vertex $(0, -1, -1)$.
        \item For $y = 0$: trace is $z = x^2$ in plane $y = 0$, vertex $(0, 0, 0)$.
        \item For $y = 1$: trace is $z = x^2 - 1$ in plane $y = 1$, vertex $(0, 1, -1)$.
    \end{itemize}
    \item In the vertical planes $x = a$, we substitute $x = a$ into $z = x^2 - y^2$:
    \[ z = a^2 - y^2. \]
    This describes a downward-opening parabola in the plane $x = a$. In this plane, coordinates are $y$ (horizontal) and $z$ (vertical). The vertex occurs at $y = 0$, giving coordinates $(x,y,z) = (a, 0, a^2)$.
    \begin{itemize}[nosep]
        \item For $x = -1$: trace is $z = 1 - y^2$ in plane $x = -1$, vertex $(-1, 0, 1)$.
        \item For $x = 0$: trace is $z = -y^2$ in plane $x = 0$, vertex $(0, 0, 0)$.
        \item For $x = 1$: trace is $z = 1 - y^2$ in plane $x = 1$, vertex $(1, 0, 1)$.
    \end{itemize}
    \item The cross-sections parallel to the $xz$-plane ($y=b$) are all upward parabolas, so moving away from the origin along the $x$-axis causes the surface to curve upward. Conversely, the cross-sections parallel to the $yz$-plane ($x=a$) are all downward parabolas, so moving away from the origin along the $y$-axis causes the surface to curve downward. At the origin $(0,0,0)$, the surface attains a local minimum along the $x$-direction and a local maximum along the $y$-direction. This combination of opposing curvatures defines a classic saddle point.
\end{enumerate}

\textbf{Answer.}
(a) Upward parabolas in planes $y = -1, 0, 1$ with equations $z = x^2-1, x^2, x^2-1$ and vertices $(0,-1,-1), (0,0,0), (0,1,-1)$.
(b) Downward parabolas in planes $x = -1, 0, 1$ with equations $z = 1-y^2, -y^2, 1-y^2$ and vertices $(-1,0,1), (0,0,0), (1,0,1)$.
(c) Orthogonal upward and downward parabolic curvatures identify the origin as a saddle point on a hyperbolic paraboloid.

\subsection{Problem Type PT-05: Sketch and Analyze Multi-Curve Families of Vertical Cross-Sections}
\textbf{Question reference.} See PT-05 in \texttt{Study-Notes.tex}.

\textbf{Worked solution.}
\begin{enumerate}[label=(\alph*),leftmargin=*]
    \item Holding $x$ constant gives vertical planes parallel to the $yz$-plane. Coordinates: horizontal $y$, vertical $z$.
    Substituting $x = -1, 0, 1$ into $f(x,y) = y^3 + xy$:
    \begin{itemize}[nosep]
        \item $x = -1$: $z = y^3 - y = y(y-1)(y+1)$. Roots at $y = -1, 0, 1$. Derivative $\frac{dz}{dy} = 3y^2 - 1 = 0 \implies y = \pm 1/\sqrt{3}$. Local maximum at $(-1/\sqrt{3}, 2/(3\sqrt{3}))$; local minimum at $(1/\sqrt{3}, -2/(3\sqrt{3}))$.
        \item $x = 0$: $z = y^3$. Root at $y=0$. Strictly increasing cubic curve with inflection point at the origin.
        \item $x = 1$: $z = y^3 + y = y(y^2+1)$. Single real root at $y=0$. Derivative $\frac{dz}{dy} = 3y^2+1 \ge 1 > 0$, strictly increasing with no turning points.
    \end{itemize}
    \item Holding $y$ constant gives vertical planes parallel to the $xz$-plane. Coordinates: horizontal $x$, vertical $z$.
    Substituting $y = -1, 0, 1$ into $f(x,y) = y^3 + xy$:
    \begin{itemize}[nosep]
        \item $y = -1$: $z = (-1)^3 + x(-1) = -1 - x$. Straight line with slope $-1$ and $z$-intercept $-1$.
        \item $y = 0$: $z = 0^3 + x(0) = 0$. Horizontal line along the $x$-axis ($z=0$).
        \item $y = 1$: $z = 1^3 + x(1) = 1 + x$. Straight line with slope $+1$ and $z$-intercept $+1$.
    \end{itemize}
    \item Along slices of constant $y$, the height varies linearly with $x$ with slope $\frac{\partial z}{\partial x} = y$. When $y < 0$, the slices slope downward; when $y = 0$, the slice is flat; when $y > 0$, the slices slope upward. Thus, moving along the $y$-axis causes the surface to twist continuously from negative slope to positive slope.
\end{enumerate}

\textbf{Answer.}
(a) On $yz$-axes: cubics $z = y^3-y$ (roots $-1,0,1$), $z = y^3$ (inflection at origin), and $z = y^3+y$ (strictly increasing).
(b) On $xz$-axes: straight lines $z = -1-x$ (slope $-1$), $z = 0$ (slope $0$), and $z = 1+x$ (slope $+1$).
(c) The slope of the $y$-fixed slices is $m = y$, causing the surface to twist about the origin.

\subsection{Problem Type PT-06: Identify and Sketch Cylinders from Missing Variables in $\R^3$}
\textbf{Question reference.} See PT-06 in \texttt{Study-Notes.tex}.

\textbf{Worked solution.}
\begin{enumerate}[label=(\alph*),leftmargin=*]
    \item In $\R^2$, $x^2 + y^2 = 1$ is a one-dimensional curve: the unit circle centered at the origin.
    In $\R^3$, the variable $z$ is absent (free). Any point $(x,y,z)$ satisfying $x^2+y^2=1$ is a solution regardless of $z$. Translating the unit circle parallel to the $z$-axis generates a right circular cylinder of radius $1$ centered on the $z$-axis. At any point on the unit circle (e.g., $(1,0)$), all points $(1,0,z)$ for $z \in \R$ lie on the surface; thus, vertical lines intersect it at infinitely many points, failing the vertical-line test for $z=f(x,y)$.
    \item In $\R^3$, $z = x^2$ lacks the variable $y$. In the coordinate plane $y=0$ ($xz$-plane), the generating curve is the standard parabola $z = x^2$. Translating this parabola along the $y$-axis produces a parabolic cylinder. For every pair $(x,y) \in \R^2$, the height $z = x^2$ is uniquely determined. Hence, $z = x^2$ \textbf{is} the graph of a valid function $z = f(x,y)$ on $\R^2$.
    \item In $\R^3$, $y = x^2$ lacks the variable $z$. In the coordinate plane $z=0$ ($xy$-plane), the generating curve is the parabola $y = x^2$. Translating this parabola parallel to the $z$-axis produces a parabolic cylinder extruded vertically. Over any point $(x_0, x_0^2)$ on the parabola, $z$ can take any real value. Thus, it completely fails the vertical-line test and cannot be written as $z = f(x,y)$.
\end{enumerate}

\textbf{Answer.}
(a) Circle in $\R^2$; vertical circular cylinder of radius $1$ around the $z$-axis in $\R^3$ (fails $z$-function test).
(b) Parabolic cylinder parallel to the $y$-axis with generating curve $z=x^2$; represents a valid function $z=f(x,y)$ on $\R^2$.
(c) Parabolic cylinder parallel to the $z$-axis with generating curve $y=x^2$; fails the $z$-function test.

\subsection{Problem Type PT-07: Distinguish Vertical Cross-Sections from Horizontal Level Curves}
\textbf{Question reference.} See PT-07 in \texttt{Study-Notes.tex}.

\textbf{Worked solution.}
\begin{enumerate}[label=(\alph*),leftmargin=*]
    \item For $f(x,y) = x^2 + y^2$, setting $x = 1$ gives the cross-section:
    \[ x = 1, \quad z = 1 + y^2. \]
    This curve lies in the vertical plane $x = 1$. The variable $x$ is held fixed at $1$, while $y$ varies over $\R$ and $z$ varies over $[1,\infty)$. It is plotted on $yz$-coordinate axes as an upward-opening parabola with vertex $(y,z) = (0,1)$.
    \item Setting $z = 4$ gives the contour curve:
    \[ z = 4, \quad x^2 + y^2 = 4. \]
    In 3D space, this curve lies in the horizontal plane $z = 4$. On the contour diagram, it is projected vertically into the $xy$-domain plane as a circle of radius $r = \sqrt{4} = 2$ centered at the origin $(0,0)$. The height $z$ is strictly fixed at $4$ along the entire curve.
    \item Two fundamental mathematical differences:
    \begin{itemize}[nosep]
        \item \emph{Fixed quantity and spatial plane:} A vertical cross-section fixes an \textbf{input} variable ($x$ or $y$) and lies in a vertical plane perpendicular to the domain; a contour fixes the \textbf{output} variable ($z$) and lies in a horizontal plane parallel to the domain.
        \item \emph{Height behavior:} Along a vertical cross-section, the output height $z$ varies; along a contour curve, the output height $z$ is strictly constant.
    \end{itemize}
\end{enumerate}

\textbf{Answer.}
(a) Vertical parabola $z = 1+y^2$ in the plane $x=1$, viewed on $yz$-axes.
(b) Horizontal circle $x^2+y^2=4$ at height $z=4$, projected onto $xy$-axes as a radius-$2$ circle.
(c) Cross-sections fix inputs and vary in height; contours fix output height and map domain curves where $z$ is constant.

\subsection{Problem Type PT-08: Find Algebraic Contours and Interpret Spacing and Steepness}
\textbf{Question reference.} See PT-08 in \texttt{Study-Notes.tex}.

\textbf{Worked solution.}
\begin{enumerate}[label=(\alph*),leftmargin=*]
    \item For $f(x,y) = x^2 + y^2$:
    Setting $x^2 + y^2 = c$ requires $c \ge 0$ since a sum of squares cannot be negative.
    \begin{itemize}[nosep]
        \item For $c < 0$: the level set is the empty set $\emptyset$.
        \item For $c = 0$: $x^2+y^2=0 \implies (x,y)=(0,0)$ (a single degenerate point, radius $0$).
        \item For $c > 0$: circles centered at $(0,0)$ of radius $r = \sqrt{c}$.
        \item Levels: $c=0 \implies r = 0$; $c=2 \implies r = \sqrt{2} \approx 1.414$; $c=4 \implies r = 2$; $c=6 \implies r = \sqrt{6} \approx 2.449$; $c=8 \implies r = \sqrt{8} = 2\sqrt{2} \approx 2.828$.
    \end{itemize}
    \item For $g(x,y) = \sqrt{x^2 + y^2}$:
    Setting $\sqrt{x^2+y^2} = c$ requires $c \ge 0$ since principal square roots are non-negative.
    \begin{itemize}[nosep]
        \item For $c < 0$: empty set $\emptyset$.
        \item For $c = 0$: $(x,y)=(0,0)$, radius $r = 0$.
        \item For $c > 0$: circles of radius $r = c$.
        \item Levels: $c=0 \implies r = 0$; $c=1 \implies r = 1$; $c=2 \implies r = 2$; $c=3 \implies r = 3$.
    \end{itemize}
    \item Radial spacing analysis:
    \begin{itemize}[nosep]
        \item For $f(x,y)$: output increment is $\Delta c = 2$. Successive radial gaps are:
        $\Delta r_1 = \sqrt{2} - 0 \approx 1.414$,
        $\Delta r_2 = 2 - \sqrt{2} \approx 0.586$,
        $\Delta r_3 = \sqrt{6} - 2 \approx 0.449$,
        $\Delta r_4 = \sqrt{8} - \sqrt{6} \approx 0.379$.
        The radial distance between adjacent contours shrinks outward ($\Delta r \to 0$). Because equal rises in height occur over progressively smaller horizontal distances, the paraboloid becomes steeper outward.
        \item For $g(x,y)$: output increment is $\Delta c = 1$. Successive radial gaps are:
        $\Delta r = 1 - 0 = 2 - 1 = 3 - 2 = 1.0$ (constant).
        The contours are equally spaced, which means that height changes at a strictly constant rate with respect to horizontal distance; the graph of $g$ is a circular cone with constant slope.
    \end{itemize}
\end{enumerate}

\textbf{Answer.}
(a) Concentric circles with radii $0, \sqrt{2}, 2, \sqrt{6}, 2\sqrt{2}$ for $c = 0, 2, 4, 6, 8$; empty for $c<0$.
(b) Concentric circles with radii $0, 1, 2, 3$ for $c = 0, 1, 2, 3$; empty for $c<0$.
(c) For $f$, shrinking radial gaps demonstrate increasing steepness (parabolic bowl); for $g$, uniform radial gaps demonstrate constant slope (cone).

\subsection{Problem Type PT-09: Plot and Interpret Contours of Linear Functions}
\textbf{Question reference.} See PT-09 in \texttt{Study-Notes.tex}.

\textbf{Worked solution.}
\begin{enumerate}[label=(\alph*),leftmargin=*]
    \item Set $z = 2y - x = c$ and solve for $y$:
    \[ 2y = x + c \implies y = \frac{1}{2}x + \frac{c}{2}. \]
    Substituting $c = 0, 2, 4, 6$:
    \begin{itemize}[nosep]
        \item $z = 0: \quad y = \frac{1}{2}x$
        \item $z = 2: \quad y = \frac{1}{2}x + 1$
        \item $z = 4: \quad y = \frac{1}{2}x + 2$
        \item $z = 6: \quad y = \frac{1}{2}x + 3$
    \end{itemize}
    \item All lines have common slope $m = 1/2$. The $y$-intercepts (where $x=0$) are $(0,0)$, $(0,1)$, $(0,2)$, and $(0,3)$, respectively.
    \item For $z = -x + 2y$, the change in height is $\Delta z = -\Delta x + 2\Delta y$.
    To increase $z$, we need $-\Delta x + 2\Delta y > 0$. Moving in the direction of the vector $\mathbf{v} = (-1, 2)$ yields:
    \[ \Delta z = -(-1) + 2(2) = 1 + 4 = 5 > 0. \]
    Thus, $z$ increases most rapidly toward the upper-left (increasing $y$, decreasing $x$).
    Moving along the line's direction vector $\mathbf{u} = (2, 1)$ yields:
    \[ \Delta z = -(2) + 2(1) = 0. \]
    This confirms that moving along a contour line produces zero change in $z$.
\end{enumerate}

\textbf{Answer.}
(a) Four parallel lines: $y = \frac{1}{2}x$, $y = \frac{1}{2}x+1$, $y = \frac{1}{2}x+2$, $y = \frac{1}{2}x+3$.
(b) Common slope is $1/2$; $y$-intercepts are $0, 1, 2, 3$.
(c) Vector direction of increasing $z$ is $(-1, 2)$; along the vector $(2, 1)$ the change $\Delta z = 0$.

\subsection{Problem Type PT-10: Read a Labelled Contour Map and Bracket / Interpolate Values}
\textbf{Question reference.} See PT-10 in \texttt{Study-Notes.tex}.

\textbf{Worked solution.}
\begin{enumerate}[label=(\alph*),leftmargin=*]
    \item Point $A = (1.5, 0)$ lies on the positive $x$-axis at distance $r = 1.5$ from the origin.
    The surrounding drawn contours for $q(x,y) = 12 - x^2 - y^2$ are:
    \begin{itemize}[nosep]
        \item $q = 10 \implies x^2+y^2 = 12-10 = 2 \implies r = \sqrt{2} \approx 1.414$.
        \item $q = 8 \implies x^2+y^2 = 12-8 = 4 \implies r = 2.0$.
    \end{itemize}
    Since $1.414 < 1.5 < 2.0$, point $A$ lies strictly between the $q=10$ and $q=8$ contours.
    Therefore, the height is rigorously bounded by:
    \[ 8 < q(1.5,0) < 10. \]
    \item Assumption: Assume that $q$ varies approximately linearly with respect to radial distance $r$ across the narrow interval $[\sqrt{2}, 2]$.
    The fraction of the distance from $r = \sqrt{2}$ to $r = 2$ is:
    \[ t = \frac{1.5 - \sqrt{2}}{2 - \sqrt{2}} \approx \frac{1.5 - 1.4142}{2 - 1.4142} = \frac{0.0858}{0.5858} \approx 0.1465. \]
    Since $q$ decreases from $10$ to $8$ as $r$ increases, the linearly interpolated value is:
    \[ q_{\text{est}} \approx 10 - t(10 - 8) = 10 - 0.1465(2) = 10 - 0.293 = 9.707 \approx 9.71. \]
    \item Exact evaluation:
    \[ q(1.5, 0) = 12 - (1.5)^2 - 0^2 = 12 - 2.25 = 9.75. \]
    The slight difference ($9.71$ vs $9.75$) occurs because the true profile $q(r) = 12 - r^2$ is concave downward ($q''(r) = -2 < 0$), so the secant line used in linear interpolation lies strictly below the true quadratic curve.
\end{enumerate}

\textbf{Answer.}
(a) Enclosing levels are $10$ and $8$; bounding inequality is $8 < q(1.5,0) < 10$.
(b) Linear radial interpolation yields $q(1.5,0) \approx 9.71$.
(c) Exact formula evaluation gives $q(1.5,0) = 9.75$. The interpolation is an estimate due to downward curvature.

\subsection{Problem Type PT-11: Construct and Interpolate Contours from Numerical Data Tables}
\textbf{Question reference.} See PT-11 in \texttt{Study-Notes.tex}.

\textbf{Worked solution.}
\begin{enumerate}[label=(\alph*),leftmargin=*]
    \item We locate where $f(x,y) = 6$ crosses each horizontal row of the table:
    \begin{itemize}[nosep]
        \item \emph{Row $y = 0$:} Entries are $f(0,0)=12$, $f(2,0)=8$, $f(4,0)=4$. The value $6$ lies between $x=2$ (value $8$) and $x=4$ (value $4$). By linear interpolation:
        \[ x = 2 + \frac{6 - 8}{4 - 8}(4 - 2) = 2 + \frac{-2}{-4}(2) = 2 + 1 = 3. \implies (3, 0). \]
        \item \emph{Row $y = 2$:} Entries are $f(0,2)=10$, $f(2,2)=6$, $f(4,2)=2$. The value $6$ occurs exactly at the grid point $x = 2$. $\implies (2, 2)$.
        \item \emph{Row $y = 4$:} Entries are $f(0,4)=8$, $f(2,4)=4$, $f(4,4)=0$. The value $6$ lies between $x=0$ (value $8$) and $x=2$ (value $4$). By linear interpolation:
        \[ x = 0 + \frac{6 - 8}{4 - 8}(2 - 0) = 0 + \frac{-2}{-4}(2) = 1. \implies (1, 4). \]
    \end{itemize}
    \item The points $(3,0)$, $(2,2)$, and $(1,4)$ are plotted in the $xy$-plane. They lie along a single straight line segment connecting the bottom edge $(3,0)$ to the top edge $(1,4)$ of the square domain $[0,4] \times [0,4]$.
    \item The slope of the line in the $xy$-plane is:
    \[ m = \frac{2 - 0}{2 - 3} = \frac{2}{-1} = -2. \]
    Using point-slope form with $(3,0)$:
    \[ y - 0 = -2(x - 3) \implies y = -2x + 6 \implies 2x + y = 6. \]
    \emph{Modeling limit:} This equation represents a consistent linear model that matches the interpolated points on this discrete grid. However, a finite table cannot prove that $f(x,y) = 12 - 2x - y$ everywhere, as nonlinear functions could share identical values at the sampled grid points.
\end{enumerate}

\textbf{Answer.}
(a) Interpolated points are $(3,0)$, $(2,2)$, and $(1,4)$.
(b) The contour is a straight line segment across the sampled square.
(c) The fitted contour equation is $2x + y = 6$; it is an empirical linear approximation valid within $[0,4]^2$.

\subsection{Problem Type PT-12: Analyze Saddle Contours, Hyperbolic Branches, and Intersecting Zero Levels}
\textbf{Question reference.} See PT-12 in \texttt{Study-Notes.tex}.

\textbf{Worked solution.}
\begin{enumerate}[label=(\alph*),leftmargin=*]
    \item From the $7 \times 7$ table for $f(x,y) = x^2 - y^2$, the entries equal to $0$ occur at:
    \[ (-3,-3), (-2,-2), (-1,-1), (0,0), (1,1), (2,2), (3,3) \implies y = x, \]
    \[ (-3,3), (-2,2), (-1,1), (0,0), (1,-1), (2,-2), (3,-3) \implies y = -x. \]
    Algebraically, $x^2 - y^2 = 0 \iff (x-y)(x+y) = 0 \iff y = x \text{ or } y = -x$. These are two perpendicular straight lines intersecting at the origin.
    \item Conic classification by level $c$:
    \begin{itemize}[nosep]
        \item For $c = 4$: $x^2 - y^2 = 4$, a hyperbola opening along the $x$-axis (left and right) with vertices at $(\pm 2, 0)$ (two disconnected branches).
        \item For $c = 2$: $x^2 - y^2 = 2$, a hyperbola opening along the $x$-axis with vertices at $(\pm\sqrt{2}, 0)$ (two disconnected branches).
        \item For $c = 0$: the two intersecting lines $y = \pm x$ dividing the plane into four quadrants (one connected geometric set).
        \item For $c = -2$: $x^2 - y^2 = -2 \iff y^2 - x^2 = 2$, a hyperbola opening along the $y$-axis (top and bottom) with vertices at $(0, \pm\sqrt{2})$ (two disconnected branches).
        \item For $c = -4$: $x^2 - y^2 = -4 \iff y^2 - x^2 = 4$, a hyperbola opening along the $y$-axis with vertices at $(0, \pm 2)$ (two disconnected branches).
    \end{itemize}
    \item The non-intersecting rule states that two contour lines corresponding to \emph{different} heights ($c_1 \ne c_2$) can never intersect. At the origin, both lines $y = x$ and $y = -x$ belong to the \textbf{same} level set $c = 0$. The function evaluates to $f(0,0) = 0$, so single-valuedness is fully preserved.
\end{enumerate}

\textbf{Answer.}
(a) Zero entries lie along the two diagonal lines $y = x$ and $y = -x$.
(b) For $c = 2, 4$, horizontal hyperbolas (vertices on $x$-axis); for $c = -2, -4$, vertical hyperbolas (vertices on $y$-axis); for $c = 0$, two intersecting diagonals.
(c) The intersection at $(0,0)$ involves curves of the same value ($c=0$), which does not violate the function definition.

\subsection{Problem Type PT-13: Model a Finite-Domain Surface, Bounded Contours, and 3D Distances}
\textbf{Question reference.} See PT-13 in \texttt{Study-Notes.tex}.

\textbf{Worked solution.}
\begin{enumerate}[label=(\alph*),leftmargin=*]
    \item The base spans $0 \le x \le 2$ and $0 \le y \le 3$. The ridge connects $(1,0,4)$ and $(1,3,4)$ at height $4$.
    \begin{itemize}[nosep]
        \item On $0 \le x \le 1$: height rises from $0$ at $x=0$ to $4$ at $x=1$ independent of $y$. Slope is $\frac{4-0}{1-0} = 4 \implies h(x,y) = 4x$.
        \item On $1 \le x \le 2$: height falls from $4$ at $x=1$ to $0$ at $x=2$ independent of $y$. Slope is $\frac{0-4}{2-1} = -4 \implies h(x,y) = 4 - 4(x-1) = 8 - 4x$.
    \end{itemize}
    Thus, the complete piecewise height function is:
    \[ h(x,y) = \begin{cases} 4x & 0 \le x \le 1, \\ 8 - 4x & 1 \le x \le 2, \end{cases} \quad \text{for } 0 \le y \le 3. \]
    \item Solving $h(x,y) = c$ for $0 < c < 4$:
    \[ 4x = c \implies x = \frac{c}{4}, \qquad 8 - 4x = c \implies x = 2 - \frac{c}{4}. \]
    Each contour consists of two vertical line segments parallel to the $y$-axis:
    \begin{itemize}[nosep]
        \item $c = 0$: $x = 0$ and $x = 2$ ($0 \le y \le 3$) (ground edges).
        \item $c = 1$: $x = 1/4$ and $x = 7/4$ ($0 \le y \le 3$).
        \item $c = 2$: $x = 1/2$ and $x = 3/2$ ($0 \le y \le 3$).
        \item $c = 3$: $x = 3/4$ and $x = 5/4$ ($0 \le y \le 3$).
        \item $c = 4$: $x = 1$ ($0 \le y \le 3$) (ridge line).
    \end{itemize}
    \item Shoes are at $S = (0, 1, 0)$. Bug is at $(1/2, 1, h(1/2,1))$.
    Height of bug: $h(1/2, 1) = 4(1/2) = 2$ m. Thus, bug coordinates are $B = (1/2, 1, 2)$.
    The 3D Euclidean distance is:
    \[ d(S, B) = \sqrt{\left(\frac{1}{2} - 0\right)^2 + (1 - 1)^2 + (2 - 0)^2} = \sqrt{\frac{1}{4} + 0 + 4} = \sqrt{\frac{17}{4}} = \frac{\sqrt{17}}{2} \text{ m} \approx 2.062 \text{ m}. \]
\end{enumerate}

\textbf{Answer.}
(a) $h(x,y) = \begin{cases} 4x & 0 \le x \le 1 \\ 8 - 4x & 1 \le x \le 2 \end{cases}$ with $0 \le y \le 3$.
(b) Pairs of vertical segments at $x = c/4$ and $x = 2 - c/4$ with $0 \le y \le 3$ for $c \in (0,4)$; single segment $x=1$ for $c=4$.
(c) Distance between shoes and bug is $\frac{\sqrt{17}}{2}$ m ($\approx 2.062$ m).

\subsection{Problem Type PT-14: Analyze Multivariable Economic Models and Trade-offs along Contours}
\textbf{Question reference.} See PT-14 in \texttt{Study-Notes.tex}.

\textbf{Worked solution.}
\begin{enumerate}[label=(\alph*),leftmargin=*]
    \item Set $P(N,V) = \sqrt{NV} = 4$.
    Squaring both sides yields:
    \[ NV = 16 \implies V = \frac{16}{N} \quad (N > 0). \]
    This is a rectangular hyperbola in the first quadrant of the $NV$-plane.
    \item If $N = 0$, then $P(0,V) = \sqrt{0 \cdot V} = 0 \ne 4$. Similarly, if $V = 0$, $P(N,0) = \sqrt{N \cdot 0} = 0 \ne 4$.
    Economically, both labor and capital are indispensable inputs; without either input, production drops to zero. As $N \to 0^+$, $V = 16/N \to \infty$; as $N \to \infty$, $V \to 0^+$. Thus, the isoquant is asymptotic to both axes but never touches or intersects them.
    \item Let the initial production point be $(N_0, V_0)$ with $N_0 V_0 = 16$.
    If labor is doubled to $N_{\text{new}} = 2N_0$, the new capital $V_{\text{new}}$ must satisfy:
    \[ N_{\text{new}} V_{\text{new}} = 16 \implies (2N_0) V_{\text{new}} = 16 \implies V_{\text{new}} = \frac{16}{2N_0} = \frac{1}{2}\left(\frac{16}{N_0}\right) = \frac{1}{2}V_0. \]
    Therefore, capital must be halved to maintain the exact same production output of $4$.
\end{enumerate}

\textbf{Answer.}
(a) The isoquant equation is $V = 16/N$ for $N > 0$.
(b) Intersecting an axis requires $N=0$ or $V=0$, yielding $P=0 \ne 4$; hence, contours never meet the axes.
(c) Capital must be halved ($V \mapsto V/2$).

\subsection{Problem Type PT-15: Describe and Visualize Level Surfaces of Three-Variable Functions}
\textbf{Question reference.} See PT-15 in \texttt{Study-Notes.tex}.

\textbf{Worked solution.}
\begin{enumerate}[label=(\alph*),leftmargin=*]
    \item For a three-variable function $w = T(x,y,z)$, a graph consists of points $(x,y,z,w) \in \R^4$, which requires four mutually perpendicular spatial dimensions and cannot be drawn in 3D space.
    A level surface $T(x,y,z) = c$ fixes the output value to a constant $c$. It forms a two-dimensional geometric surface embedded in physical 3D space where the temperature is uniform.
    \item For $T(x,y,z) = \frac{1}{4}(x^2 + y^2 + z^2) = c \implies x^2 + y^2 + z^2 = 4c$.
    \begin{itemize}[nosep]
        \item For $T = 1^\circ\text{C}$: $x^2 + y^2 + z^2 = 4(1) = 4 = 2^2$. A sphere of radius $r = 2$ m centered at $(0,0,0)$.
        \item For $T = 5^\circ\text{C}$: $x^2 + y^2 + z^2 = 4(5) = 20 = (\sqrt{20})^2$. A sphere of radius $r = 2\sqrt{5} \approx 4.47$ m centered at $(0,0,0)$.
        \item For $T = 9^\circ\text{C}$: $x^2 + y^2 + z^2 = 4(9) = 36 = 6^2$. A sphere of radius $r = 6$ m centered at $(0,0,0)$.
    \end{itemize}
    \item Level surfaces for additional functions:
    \begin{itemize}[nosep]
        \item (i) $F(x,y,z) = e^{-(x^2+y^2)} = c$:
        Requires $0 < c \le 1$. Taking natural logarithms: $-(x^2+y^2) = \ln c \implies x^2 + y^2 = -\ln c = \ln(1/c)$.
        Since $z$ is absent (free), for each $c \in (0,1)$ this is a circular cylinder of radius $\sqrt{\ln(1/c)}$ parallel to the $z$-axis. For $c = 1$, $-\ln 1 = 0$, giving the $z$-axis ($x=0, y=0$). For $c \le 0$ or $c > 1$, the level set is empty.
        \item (ii) $G(x,y,z) = z - y = c \implies z = y + c$:
        Linear equation in three variables where $x$ is free. For each $c \in \R$, this is a flat infinite plane parallel to the $x$-axis, inclined at an angle of $45^\circ$ to the $y$- and $z$-axes.
        \item (iii) $H(x,y,z) = \ln(x^2+y^2+z^2) = c \implies x^2 + y^2 + z^2 = e^c$:
        For any real $c \in \R$, $e^c > 0$. The level surface is a sphere centered at $(0,0,0)$ with radius $r = e^{c/2}$. The origin $(0,0,0)$ is excluded from the domain of $H$.
    \end{itemize}
\end{enumerate}

\textbf{Answer.}
(a) The 4D graph is not visualizable; level surfaces are 2D surfaces in 3D space where temperature is constant.
(b) Spheres centered at origin with radii $2$ m, $2\sqrt{5}$ m, and $6$ m.
(c) (i) Circular cylinders parallel to the $z$-axis for $0<c<1$; (ii) Parallel planes inclined at $45^\circ$ to the $yz$-axes; (iii) Concentric spheres of radius $e^{c/2}$.

\clearpage

% =========================================================================
% SECTION B: ACTUAL WEEKLY QUESTIONS AND SOLUTIONS
% =========================================================================
\section{Actual Weekly Questions and Solutions}\label{sec:weekly-questions}
This section contains all questions extracted verbatim from the weekly lecture slides, handwritten notes, and term test problems.

\subsection{\texttt{Week2.pdf}, Page 1 --- Describing a Plane}
\textbf{Question.} Describe the graph of $f(x,y) = 1 + x - y$. Work with cross-sections first; use them to predict the full surface. Fill the $x=a$ and $y=b$ cross-section equations and the directions of change.

\textbf{Relevant method.} PT-01: Recognize and Describe a Surface from an Equation; PT-04: Vertical Traces.

\textbf{Necessary work.}
Fix $x=a$: $z = 1 + a - y$, a straight line in the plane $x=a$ with slope $-1$ in $y$, decreasing as $y$ increases.
Fix $y=b$: $z = 1 + x - b$, a straight line in the plane $y=b$ with slope $+1$ in $x$, increasing as $x$ increases.

\textbf{Answer.}
The graph is an infinite plane sloping downward in the positive $y$-direction and upward in the positive $x$-direction, passing through $(0,0,1)$ with normal vector $(-1, 1, 1)$.

\subsection{\texttt{Week2.pdf}, Page 2 --- Common Surfaces and Ambient Space}
\textbf{Question.} For each equation, name the surface and decide whether the whole surface is the graph of a function $z = f(x,y)$:
$z = 2 - x + 3y$, $z = x^2 + y^2$, $x^2 + y^2 + z^2 = 4$, $z = x^2$.
Graph $x^2 + y^2 = 1$ first in $\R^2$, then in $\R^3$; use the missing-variable cylinder test.

\textbf{Relevant method.} PT-01: Recognize and Describe a Surface; PT-06: Cylinders.

\textbf{Necessary work.}
\begin{center}
\begin{tabular}{lll}
\toprule
Equation & Surface Type & $z = f(x,y)$? \\
\midrule
$z = 2 - x + 3y$ & Plane & Yes \\
$z = x^2 + y^2$ & Circular paraboloid & Yes \\
$x^2 + y^2 + z^2 = 4$ & Sphere (radius 2) & No \\
$z = x^2$ & Parabolic cylinder (parallel to $y$) & Yes \\
\bottomrule
\end{tabular}
\end{center}
For $x^2+y^2=1$: in $\R^2$, it is a 1D unit circle; in $\R^3$, $z$ is free, extruding the circle vertically into a circular cylinder, which fails the vertical-line test for $z=f(x,y)$.

\textbf{Answer.}
Classification as in table. $x^2+y^2=1$ is a circle in $\R^2$ and a circular cylinder parallel to $z$ in $\R^3$ (not a $z$-function).

\subsection{\texttt{Week2.pdf}, Page 3 --- Six Cross-Sections}
\textbf{Question.} Let $f(x,y) = y^3 + xy$. Find the indicated equations, then draw graphs of each family of cross-sections on the same plane: (a) fix $x = -1, 0, 1$; (b) fix $y = -1, 0, 1$.

\textbf{Relevant method.} PT-05: Multi-Curve Families of Vertical Cross-Sections.

\textbf{Necessary work.}
(a) Fixed $x$: curves on $yz$-axes.
$x = -1 \implies z = y^3 - y$;
$x = 0 \implies z = y^3$;
$x = 1 \implies z = y^3 + y$.
(b) Fixed $y$: curves on $xz$-axes.
$y = -1 \implies z = -1 - x$;
$y = 0 \implies z = 0$;
$y = 1 \implies z = 1 + x$.

\textbf{Answer.}
(a) Three cubic curves on $yz$-axes: $z = y^3-y$, $z=y^3$, $z=y^3+y$.
(b) Three straight lines on $xz$-axes: $z = -1-x$, $z=0$, $z=1+x$.
\fig{.80}{03-cross-sections.png}

\subsection{\texttt{Week2.pdf}, Pages 4--5 --- Paraboloid and Cone Contours}
\textbf{Question.} Find equations for the contours of $f(x,y) = x^2 + y^2$, draw a diagram for $c = 0, 2, 4, 6, 8$, and relate spacing to the graph. Then draw contours of $g(x,y) = \sqrt{x^2+y^2}$ for $c = 0, 1, 2, 3$ and compare the cone with the paraboloid.

\textbf{Relevant method.} PT-08: Find Algebraic Contours and Spacing.

\textbf{Necessary work.}
$f(x,y) = c \implies x^2+y^2 = c \implies r = \sqrt{c}$ for $c \ge 0$. Radii are $0, \sqrt{2}, 2, \sqrt{6}, 2\sqrt{2}$.
$g(x,y) = c \implies r = c$ for $c \ge 0$. Radii are $0, 1, 2, 3$.

\textbf{Answer.}
Paraboloid contours are circles with radii $\sqrt{c}$ whose spacing shrinks outward, reflecting a bowl that steepens with distance. Cone contours are circles with radii $c$ with uniform spacing, reflecting constant slope.
\fig{.82}{04-radial-contours.png}

\subsection{\texttt{Week2.pdf}, Page 6 --- Contour Exercise for $z = 2y - x$}
\textbf{Question.} Make a contour plot for $z = 2y - x$ with at least three different contours. Label each contour with its value of $z$. The worksheet lists $c = 0, 2, 4, 6$. Which direction in the $xy$-plane gives increasing values of $z$?

\textbf{Relevant method.} PT-09: Contours of Linear Functions.

\textbf{Necessary work.}
$2y - x = c \implies y = \frac{1}{2}x + \frac{c}{2}$. Slopes are $1/2$. Intercepts are $(0, c/2)$. Vector $(-1, 2)$ gives $\Delta z = -(-1) + 2(2) = 5 > 0$.

\textbf{Answer.}
Parallel lines $y = \frac{1}{2}x + \frac{c}{2}$ for $c = 0, 2, 4, 6$. Height increases toward the upper-left (direction of $(-1, 2)$).
\fig{.80}{05-plane-saddle-contours.png}

\subsection{\texttt{Week2.pdf}, Pages 7--8 --- Numerical Saddle Table}
\textbf{Question.} Relate the table for $f(x,y) = x^2 - y^2$ to its contour diagram. (1) Locate zero entries and algebraically solve $x^2 - y^2 = 0$. (2) Mark positive and negative regions and symmetries. On the grid draw and label $z = -4, -2, 0, 2, 4$. Use the table to decide where each branch lies; relate the contours to the saddle graph.

\textbf{Relevant method.} PT-12: Analyze Saddle Contours and Tables.

\textbf{Necessary work.}
Zero entries satisfy $x^2-y^2=0 \implies y = \pm x$. $f > 0$ when $|x| > |y|$; $f < 0$ when $|y| > |x|$. For $c > 0$, hyperbolas open along $x$; for $c < 0$, hyperbolas open along $y$.

\textbf{Answer.}
Zero level is the pair of intersecting lines $y = \pm x$. Positive levels ($2, 4$) are hyperbolas opening left/right; negative levels ($-2, -4$) are hyperbolas opening top/bottom.

\subsection{\texttt{Week2.pdf}, Pages 9--10 --- Fall 2025 Term-Test Tent}
\textbf{Question.} The four corners of the tent are $A=(0,0,0)$, $B=(2,0,0)$, $C=(2,3,0)$, $D=(0,3,0)$, and the ridge endpoints are $P=(1,0,4)$, $Q=(1,3,4)$. Let $h(x,y)$ be the height, in metres, of the tent at position $(x,y)$; the ground has height $0$. Find a formula for the roof height and for $0 < c < 4$ solve $h(x,y) = c$.
\begin{enumerate}[label=(\alph*),nosep]
    \item Sketch at least three contours of $h$, with heights labelled.
    \item Shoes are at $(0,1,0)$ and a bug is at $(1/2, 1, h(1/2,1))$; find their Euclidean distance.
\end{enumerate}

\textbf{Relevant method.} PT-13: Model Finite-Domain Roof and Distances.

\textbf{Necessary work.}
Piecewise formula: $h(x,y) = 4x$ on $0 \le x \le 1$, $h(x,y) = 8-4x$ on $1 \le x \le 2$, with $0 \le y \le 3$.
Levels: $x = c/4$ and $x = 2 - c/4$ on $0 \le y \le 3$.
Bug position: $(1/2, 1, h(1/2,1)) = (1/2, 1, 2)$. Shoes: $(0, 1, 0)$.
Distance: $d = \sqrt{(1/2)^2 + 0^2 + 2^2} = \sqrt{17}/2$ m.

\textbf{Answer.}
(a) Pairs of vertical segments $x = c/4$ and $x = 2 - c/4$ on $0 \le y \le 3$.
(b) $d = \frac{\sqrt{17}}{2}$ m.
\fig{.75}{06-tent.png}

\subsection{\texttt{sep+14.pdf}, Pages 1--4 --- Function Test, Shifts, and Saddle Slices}
\textbf{Question.} Can a vertical line parallel to the $z$-axis intersect the graph $z = x^2 + y^2$ twice? Can the sphere $x^2 + y^2 + z^2 = 1$ be seen as a graph of one function? Express it as two graphs. Describe the transformations $z = x^2+y^2 \mapsto z = (x-1)^2 + (y-2)^2$, $z = x+y \mapsto z = x+y+1$, and reflection of the upper hemisphere. For the paraboloid $z = x^2+y^2$, take a slice $y = a$. Get slices for $z = x^2-y^2$.

\textbf{Relevant method.} PT-02: Test Function and Domain; PT-03: Transformations; PT-04: Parabolic Traces.

\textbf{Necessary work.}
Paraboloid passes vertical-line test since $z$ is unique. Sphere fails at $(0,0,\pm 1)$; splits into $z = \pm\sqrt{1-x^2-y^2}$ on $x^2+y^2\le 1$. Shift by $(1,2,0)$; plane raised by $1$; hemisphere reflection is $z = -\sqrt{1-x^2-y^2}$. Paraboloid slice $y=a$ is $z = x^2+a^2$ (upward parabola). Saddle slice $y=a$ is $z = x^2-a^2$ (upward parabola); slice $x=a$ is $z = a^2-y^2$ (downward parabola).

\textbf{Answer.}
Paraboloid passes; sphere fails and splits into two hemispheres on $x^2+y^2\le 1$. Transformations translate and reflect as stated. Traces are 2D parabolas in designated vertical planes.

\subsection{\texttt{sep+16.pdf}, Pages 1--4 --- Cylinders and Downward Bowl Contours}
\textbf{Question.} For $z = x^2$, what happens when $x$ is fixed and $y$ moves? Compare $z = x^2$ with $x^2 + y^2 = 4$. For $f = 9 - x^2 - y^2$, take horizontal slices $z = a$, including $a = 8, 5$, and say what happens for $a > 9$. Interpret equal levels $9, 8, 7, 6, \dots$.

\textbf{Relevant method.} PT-06: Cylinders; PT-08: Contours and Spacing.

\textbf{Necessary work.}
On $z = x^2$, holding $x=x_0$ keeps $z=x_0^2$ constant for all $y \in \R$, forming horizontal ruling lines along $y$ (parabolic cylinder parallel to $y$). The surface $x^2+y^2=4$ is a circular cylinder parallel to $z$.
For $9 - x^2 - y^2 = a \implies x^2+y^2 = 9-a$: radius is $r = \sqrt{9-a}$. For $a=8, r=1$; for $a=5, r=2$; for $a=9, r=0$ (peak point); for $a>9$, empty set. Successive outward circles become closer, showing the bowl falls faster outward.

\textbf{Answer.}
$z=x^2$ is a parabolic cylinder parallel to $y$; $x^2+y^2=4$ is a circular cylinder parallel to $z$. Bowl contours are circles of radius $\sqrt{9-a}$; for $a>9$ empty; spacing shrinks outward indicating increasing steepness.

\subsection{\texttt{sep+18.pdf}, Pages 1--4 --- Numerical Contour Interpolation and Comparison}
\textbf{Question.} Given $f$ at $x = 0, 2, 4$ and $y = 0, 2, 4$, sketch an approximate $f = 6$ contour by estimating where it crosses the sampled rows. The note also contrasts $f = x + 2y$, $f = x^2 + y^2$, and $f = x^2 - y^2$.

\textbf{Relevant method.} PT-11: Contours from Tables; PT-12: Saddle Contours.

\textbf{Necessary work.}
Row interpolation across $y=0, 2, 4$ yields crossing points $(3,0)$, $(2,2)$, and $(1,4)$, connecting along $2x+y=6$.
Contrasts: $x+2y=c$ yields parallel straight lines of slope $-1/2$; $x^2+y^2=c$ yields concentric circles; $x^2-y^2=c$ yields hyperbolas with asymptotes $y = \pm x$.

\textbf{Answer.}
Interpolated contour segment is $2x+y=6$ on $[0,4]^2$. Level geometries are parallel lines, concentric circles, and hyperbolas, respectively.

\subsection{\texttt{MAT235H-5201-Sept-16.pdf}, Pages 2--4 --- Review and Cylinder Volume}
\textbf{Question.} Contrast one-input and two-input functions using the sample table $x \in \{1, 2, 3.14\}$ with outputs $\{1, 0, 2\}$. Describe the cylinder-volume function $V(r,h) = \pi r^2 h$. Locate a point $P = (x_0, y_0, z_0)$ in 3D Cartesian space.

\textbf{Relevant method.} PT-01: Surface Description; PT-06: Cylinders.

\textbf{Necessary work.}
A single-variable function maps $x \mapsto y$; the given table assigns single values to isolated points. The cylinder-volume function $V: [0,\infty) \times [0,\infty) \to [0,\infty)$ requires two independent inputs: radius $r$ (length) and height $h$ (length), producing volume (length$^3$). Point $P$ is located by moving $x_0$ along the $x$-axis, $y_0$ parallel to the $y$-axis, and $z_0$ parallel to the $z$-axis.

\textbf{Answer.}
One-variable functions map $\R \to \R$; $V(r,h) = \pi r^2 h$ maps $\R^2 \to \R$, requiring two inputs. Cartesian coordinates locate points uniquely via orthogonal displacements from the origin.

\subsection{\texttt{MAT235H-5201-Sept-16.pdf}, Page 7 --- Wind-Chill Table}
\textbf{Question.} Let $C = f(w,T)$ be wind chill in $^\circ\text{F}$, with wind speed $w$ in mph and temperature $T$ in $^\circ\text{F}$.
\begin{center}
\small
\begin{tabular}{c|rrrrrrrr}
\toprule
$w \backslash T$ & $35$ & $30$ & $25$ & $20$ & $15$ & $10$ & $5$ & $0$ \\
\midrule
$5$  & $31$ & $25$ & $19$ & $13$ & $7$  & $1$  & $-5$  & $-11$ \\
$10$ & $27$ & $21$ & $15$ & $9$  & $3$  & $-4$ & $-10$ & $-16$ \\
$15$ & $25$ & $19$ & $13$ & $6$  & $0$  & $-7$ & $-13$ & $-19$ \\
$20$ & $24$ & $17$ & $11$ & $4$  & $-2$ & $-9$ & $-15$ & $-22$ \\
$25$ & $23$ & $16$ & $9$  & $3$  & $-4$ & $-11$& $-17$ & $-24$ \\
\bottomrule
\end{tabular}
\end{center}
Evaluate and interpret $f(20,5)$. At actual $T = 5^\circ\text{F}$, find the wind speed giving a wind chill of $-10^\circ\text{F}$.

\textbf{Relevant method.} PT-11: Data Tables.

\textbf{Necessary work.}
Locate row $w=20$ and column $T=5$: table entry is $-15$.
Locate column $T=5$ and search vertically for entry $-10$: corresponds to row $w=10$.

\textbf{Answer.}
$f(20,5) = -15^\circ\text{F}$, meaning air at $5^\circ\text{F}$ with a $20$ mph wind feels like $-15^\circ\text{F}$. At $T = 5^\circ\text{F}$, a wind speed of $w = 10$ mph produces a wind chill of $-10^\circ\text{F}$.

\subsection{\texttt{MAT235H-5201-Sept-16.pdf}, Pages 5, 6, 9, 11 --- Coordinates, Planes, and Heater Traces}
\textbf{Question.} Plot $P=(0,0,2)$, $Q=(3,4,0)$, $R=(0,4,5)$, $S=(3,4,5)$. For $f(x,y)=x^2+y^2$ evaluate at $x,y \in \{0,1,-1\}$ and sketch $g(x,y)=3$. Which of $A=(1,-1,0), B=(0,3,4), C=(2,2,1), D=(0,-4,0)$ is nearest the $xz$-plane, and which lies on the $y$-axis? For room temperature $T = f(d,t)$, where $d$ is metres from a heater and $t$ is minutes after switching it on, describe and sketch $t$ and $d$ cross-sections qualitatively.

\textbf{Relevant method.} PT-01: Surfaces; PT-04: Cross-sections; PT-05: Trace Families.

\textbf{Necessary work.}
$P$ is on the $z$-axis; $Q$ is in the $xy$-plane; $R$ is in the $yz$-plane; $S$ is in the first octant.
$g(x,y)=3$ is a horizontal plane at height $z=3$.
Distance to $xz$-plane ($y=0$) is $|y|$: $A$ has $|-1|=1$, $B$ has $3$, $C$ has $2$, $D$ has $|-4|=4$. Thus $A$ is closest. $D=(0,-4,0)$ has $x=0, z=0$, lying on the $y$-axis.
For heater: holding $d$ fixed, $T$ increases with $t$ toward an asymptote; holding $t$ fixed, $T$ decreases as $d$ increases.

\textbf{Answer.}
Positions as stated; $g=3$ is a horizontal plane. Point $A$ is closest to the $xz$-plane; point $D$ lies on the $y$-axis. Fixed-$d$ cross-sections rise over time; fixed-$t$ cross-sections decay with distance.
\fig{.72}{12-heater-cross-sections.png}
\fig{.70}{13-coordinates-plane.png}

\subsection{\texttt{MAT235H-5201-Sept-16.pdf}, Pages 13--19 --- Cylinders, Spheres, and Hemisphere Slice}
\textbf{Question.} Sketch $x^2+y^2=4$, $x^2+z^2=9$, and $y=x^2$ in 3D. State the general equation of a sphere of radius $r$ centered at the origin. For the upper hemisphere $z = \sqrt{1-x^2-y^2}$, find the cross-section in the plane $x = 1/2$ and identify the curve.

\textbf{Relevant method.} PT-04: Cross-sections; PT-06: Cylinders.

\textbf{Necessary work.}
$x^2+y^2=4$ is a circular cylinder of radius $2$ parallel to $z$; $x^2+z^2=9$ is a circular cylinder of radius $3$ parallel to $y$; $y=x^2$ is a parabolic cylinder parallel to $z$. Sphere equation: $x^2+y^2+z^2=r^2$.
Hemisphere slice at $x=1/2$: $z = \sqrt{1 - 1/4 - y^2} = \sqrt{3/4 - y^2}$ for $|y| \le \sqrt{3}/2$. In plane $x=1/2$, this is $y^2+z^2 = 3/4$ with $z \ge 0$.

\textbf{Answer.}
Cylinders extruded along missing axes as described. Sphere is $x^2+y^2+z^2=r^2$. The slice at $x=1/2$ is an upper semicircle in the plane $x=1/2$ with radius $\sqrt{3}/2$ centered at $(1/2, 0, 0)$.
\fig{.75}{11-hemisphere-slice.png}

\subsection{\texttt{MAT235H-5201-Sept-16.pdf}, Pages 24--27; Sept-17 Pages 1--3 --- Airline Table and Planes}
\textbf{Question.} From the airline-revenue table, find $R(d,f)$ for discount tickets $d$ and full fares $f$ (revenue in thousands of dollars). Find the plane through $(4,0,0), (0,3,0), (0,0,2)$. Decide which of three given numerical tables (A, B, C) could represent linear functions. Find intercepts for $z = 2 - 2x + y$ and $z = 2 - x - 2y$.
\begin{center}
\small
\begin{tabular}{c|rrrr}
\toprule
$d \backslash f$ & $0$ & $100$ & $200$ & $300$ \\
\midrule
$0$   & $0$  & $24$ & $48$ & $72$ \\
$100$ & $8$  & $32$ & $56$ & $80$ \\
$200$ & $16$ & $40$ & $64$ & $88$ \\
$300$ & $24$ & $48$ & $72$ & $96$ \\
\bottomrule
\end{tabular}
\end{center}

\textbf{Relevant method.} PT-01: Planes; PT-11: Data Tables.

\textbf{Necessary work.}
Airline revenue: $\Delta R / \Delta d = 8/100 = 0.08$; $\Delta R / \Delta f = 24/100 = 0.24$. Formula $R(d,f) = 0.08d + 0.24f$.
Intercept plane: $x/4 + y/3 + z/2 = 1 \implies z = 2 - x/2 - 2y/3$.
Table linearity test: Table A has varying differences (nonlinear); Table B has constant $\Delta z / \Delta x = 3/2$ and $\Delta z / \Delta y = -5/3$ (linear, $z = 1 + 1.5x - 1.67y$); Table C has alternating rows (nonlinear).
Intercepts for $z = 2 - 2x + y$: $(1, 0, 0), (0, -2, 0), (0, 0, 2)$.
Intercepts for $z = 2 - x - 2y$: $(2, 0, 0), (0, 1, 0), (0, 0, 2)$.

\textbf{Answer.}
Revenue: $R(d,f) = 0.08d + 0.24f$ thousand dollars. Plane: $x/4 + y/3 + z/2 = 1$. Only Table B can represent a linear function. Intercepts are $(1,-2,2)$ and $(2,1,2)$, respectively.

\subsection{\texttt{MAT235H-5201-Sept-17.pdf}, Pages 4--5 --- Level Surfaces Preview}
\textbf{Question.} In a block of ice, $T(x,y,z) = \frac{1}{4}(x^2+y^2+z^2)$ ($^\circ\text{C}$, metres). Why is its full graph beyond ordinary 3D space? Find and explain level surfaces for $T = 1, 5, 9$. Describe level surfaces of (a) $e^{-(x^2+y^2)}$, (b) $z - y$, and (c) $\ln(x^2+y^2+z^2)$.

\textbf{Relevant method.} PT-15: Level Surfaces.

\textbf{Necessary work.}
Graph requires 4D $(x,y,z,T)$.
$T = c \implies x^2+y^2+z^2 = 4c$. For $c = 1, 5, 9$, radii are $2, 2\sqrt{5}, 6$ m.
(a) $e^{-(x^2+y^2)} = c \implies x^2+y^2 = -\ln c$, circular cylinders parallel to $z$ for $0<c<1$; $z$-axis for $c=1$.
(b) $z - y = c \implies z = y + c$, parallel planes inclined at $45^\circ$ to $y$- and $z$-axes.
(c) $\ln(x^2+y^2+z^2) = c \implies x^2+y^2+z^2 = e^c$, concentric spheres of radius $e^{c/2}$.

\textbf{Answer.}
Graph requires four dimensions. $T=1,5,9$ are concentric spheres of radii $2, 2\sqrt{5}, 6$ m. Functions (a), (b), (c) yield circular cylinders, parallel planes, and concentric spheres, respectively.
\fig{.75}{09-level-surfaces.png}

\clearpage

% =========================================================================
% SECTION C: ASSIGNED TEXTBOOK EXERCISES AND SOLUTIONS
% =========================================================================
\section{Assigned Textbook Exercises and Solutions}\label{sec:textbook-practice}
This section provides complete wording, diagrams, concise necessary work, and final answers for all 26 assigned exercises from Sections 12.2 and 12.3 of \emph{Calculus, Eighth Edition}.

\subsection{12.2, Exercise 1: Point Membership on Constant Function}
\textbf{Question.} Which of (I)--(IV) lie on the graph of the function $z = f(x,y)$, where $f(x,y) = -3$?
\[ \mathrm{I}: (1,0,1), \quad \mathrm{II}: (\sqrt{8}, 1, 3), \quad \mathrm{III}: (-3, 7, -3), \quad \mathrm{IV}: (1, 1, 1/2). \]

\textbf{Relevant method.} PT-01: Recognize and Describe a Surface from an Equation.

\textbf{Necessary work.}
The graph is $G_f = \{(x,y,z) \in \R^3 : z = -3\}$. Any point on the graph must have its third coordinate $z$ equal to $-3$. Among the choices, only III has $z = -3$.

\textbf{Answer.}
Only \textbf{III}, $(-3, 7, -3)$.

\subsection{12.2, Exercise 3: Point Membership on Reciprocal Surface}
\textbf{Question.} Which of (I)--(IV) lie on the graph of the function $z = f(x,y)$, where $f(x,y) = 1/(x^2+y^2)$?
\[ \mathrm{I}: (1,0,1), \quad \mathrm{II}: (\sqrt{8}, 1, 3), \quad \mathrm{III}: (-3, 7, -3), \quad \mathrm{IV}: (1, 1, 1/2). \]

\textbf{Relevant method.} PT-01: Recognize and Describe a Surface from an Equation.

\textbf{Necessary work.}
Evaluate $f(x,y)$ at the $(x,y)$ coordinates of each candidate point:
\begin{itemize}[nosep]
    \item I: $f(1,0) = 1/(1^2+0^2) = 1$. Matches $z=1$. $\implies$ Lies on graph.
    \item II: $f(\sqrt{8}, 1) = 1/(8+1) = 1/9 \ne 3$. $\implies$ Does not lie on graph.
    \item III: $f(-3, 7) = 1/(9+49) = 1/58 \ne -3$. $\implies$ Does not lie on graph.
    \item IV: $f(1,1) = 1/(1+1) = 1/2$. Matches $z=1/2$. $\implies$ Lies on graph.
\end{itemize}

\textbf{Answer.}
\textbf{I and IV}.

\subsection{12.2, Exercise 5: Match Functions with Surfaces}
\textbf{Question.} Without a calculator or computer, match the functions with their graphs in Figure 12.27.
\[
\begin{array}{ll}
\text{(a) } z = 2 + x^2 + y^2, & \text{(b) } z = 2 - x^2 - y^2, \\
\text{(c) } z = 2(x^2 + y^2),   & \text{(d) } z = 2 + 2x - y, \\
\text{(e) } z = 2.              &
\end{array}
\]
\fig{.55}{source-12-2-5.png}

\textbf{Relevant method.} PT-01: Surface Description; PT-03: Transformations.

\textbf{Necessary work.}
\begin{itemize}[nosep]
    \item (a) $z = 2 + x^2+y^2$ is an upward paraboloid with vertex $(0,0,2) \implies$ Graph IV.
    \item (b) $z = 2 - (x^2+y^2)$ is a downward paraboloid with vertex $(0,0,2) \implies$ Graph II.
    \item (c) $z = 2(x^2+y^2)$ is an upward paraboloid with vertex at origin $(0,0,0) \implies$ Graph I.
    \item (d) $z = 2 + 2x - y$ is an inclined plane with non-zero slopes $\implies$ Graph V.
    \item (e) $z = 2$ is a horizontal plane at height $2 \implies$ Graph III.
\end{itemize}

\textbf{Answer.}
(a) $\to$ IV; (b) $\to$ II; (c) $\to$ I; (d) $\to$ V; (e) $\to$ III.

\subsection{12.2, Exercise 7: Read Cross-Sections from a Surface Graph}
\textbf{Question.} Figure 12.29 shows the graph of $z = f(x,y)$.
\begin{enumerate}[label=(\alph*),nosep]
    \item Suppose $y$ is fixed and positive. Does $z$ increase or decrease as $x$ increases? Graph $z$ against $x$.
    \item Suppose $x$ is fixed and positive. Does $z$ increase or decrease as $y$ increases? Graph $z$ against $y$.
\end{enumerate}
\fig{.32}{source-12-2-7.png}

\textbf{Relevant method.} PT-05: Multi-Curve Families of Vertical Cross-Sections.

\textbf{Necessary work.}
(a) Tracking parallel to the positive $x$-axis at fixed positive $y$, the surface slopes downward. Hence $z$ decreases as $x$ increases.
(b) Tracking parallel to the positive $y$-axis at fixed positive $x$, the surface curls upward. Hence $z$ increases as $y$ increases.

\textbf{Answer.}
(a) $z$ decreases as $x$ increases; the cross-section bends downward.
(b) $z$ increases as $y$ increases; the cross-section bends upward.
\fig{.75}{answer-12-2-7.png}

\subsection{12.2, Exercise 9: Sketch and Describe Sphere}
\textbf{Question.} Sketch a graph of the surface and briefly describe it in words: $x^2 + y^2 + z^2 = 9$.

\textbf{Relevant method.} PT-01: Surface Description; PT-02: Function Test.

\textbf{Necessary work.}
Standard equation of a sphere $(x-x_0)^2+(y-y_0)^2+(z-z_0)^2=r^2$ with center $(0,0,0)$ and $r^2 = 9 \implies r = 3$. Axis intercepts are $(\pm 3, 0, 0)$, $(0, \pm 3, 0)$, and $(0, 0, \pm 3)$. Fails vertical line test at $(0,0,\pm 3)$.

\textbf{Answer.}
A sphere of radius $3$ centered at the origin $(0,0,0)$. Coordinate traces are circles of radius $3$. The surface fails the vertical-line test and is not the graph of a single function. (See multi-panel sketch below Exercise 15).

\subsection{12.2, Exercise 11: Sketch and Describe Downward Paraboloid}
\textbf{Question.} Sketch a graph of the surface and briefly describe it in words: $z = 5 - x^2 - y^2$.

\textbf{Relevant method.} PT-01: Surface Description; PT-03: Transformations.

\textbf{Necessary work.}
Rewrite as $z = 5 - (x^2+y^2)$. Vertex is at $(0,0,5)$. Cross-sections with $x=a$ or $y=b$ are downward parabolas. Horizontal cross-sections with $z=c$ ($c \le 5$) are circles $x^2+y^2 = 5-c$. The intersection with the $xy$-plane ($z=0$) is a circle of radius $\sqrt{5}$.

\textbf{Answer.}
A downward-opening circular paraboloid with vertex at $(0,0,5)$ and vertical axis of symmetry along the $z$-axis, intersecting the $xy$-plane in the circle $x^2+y^2=5$.

\subsection{12.2, Exercise 13: Sketch and Describe Plane}
\textbf{Question.} Sketch a graph of the surface and briefly describe it in words: $2x + 4y + 3z = 12$.

\textbf{Relevant method.} PT-01: Surface Description; PT-09: Linear Functions.

\textbf{Necessary work.}
Divide by $12$ to obtain intercept form: $\frac{x}{6} + \frac{y}{3} + \frac{z}{4} = 1$.
Intercepts: $(6,0,0)$, $(0,3,0)$, $(0,0,4)$.
Solving for $z$: $z = 4 - \frac{2}{3}x - \frac{4}{3}y$.

\textbf{Answer.}
An infinite plane with axis intercepts $(6,0,0)$, $(0,3,0)$, and $(0,0,4)$, sloping downward in both positive $x$ and positive $y$ directions.

\subsection{12.2, Exercise 15: Sketch and Describe Cylinder}
\textbf{Question.} Sketch a graph of the surface and briefly describe it in words: $x^2 + z^2 = 4$.

\textbf{Relevant method.} PT-06: Cylinders from Missing Variables.

\textbf{Necessary work.}
The variable $y$ is missing from the equation, meaning $y$ is free. In the $xz$-plane ($y=0$), the curve is a circle of radius $2$ centered at the origin: $x^2+z^2=4$. Translating this circle parallel to the $y$-axis generates a circular cylinder.

\textbf{Answer.}
A right circular cylinder of radius $2$ with central axis along the $y$-axis, extending infinitely in both positive and negative $y$ directions.
\fig{.90}{answer-12-2-surfaces.png}

\subsection{12.2, Exercise 17: Equation of Translated Sphere}
\textbf{Question.} Find the equation of the surface: a sphere of radius $3$ centered at $(0, \sqrt{7}, 0)$.

\textbf{Relevant method.} PT-01: Surface Description; PT-03: Transformations.

\textbf{Necessary work.}
Standard sphere formula: $(x - x_0)^2 + (y - y_0)^2 + (z - z_0)^2 = r^2$.
Substitute $(x_0, y_0, z_0) = (0, \sqrt{7}, 0)$ and $r = 3$:
\[ (x - 0)^2 + (y - \sqrt{7})^2 + (z - 0)^2 = 3^2. \]

\textbf{Answer.}
$\boxed{x^2 + (y - \sqrt{7})^2 + z^2 = 9}$.

\subsection{12.2, Exercise 21: Drug Concentration Traces}
\textbf{Question.} The concentration $C$, in mg/liter, of a drug in the blood is a function of $x$, the amount in mg of the drug given, and $t$, the time in hours since injection. For $0 \le x \le 4$ and $t \ge 0$,
\[ C = f(x,t) = t e^{-t(5-x)}. \]
Graph the following single-variable functions and explain their significance in terms of drug concentration: (a) $f(4,t)$; (b) $f(x,1)$.

\textbf{Relevant method.} PT-04: Cross-sections; PT-05: Multi-Curve Families.

\textbf{Necessary work.}
(a) Fix $x = 4$: $f(4,t) = t e^{-t(5-4)} = t e^{-t}$ for $t \ge 0$.
Derivative: $\frac{d}{dt}(t e^{-t}) = e^{-t}(1-t) = 0 \implies t = 1$. Peak concentration is $1/e \approx 0.368$ mg/liter at $t = 1$ hr. As $t \to \infty$, $C \to 0$.
(b) Fix $t = 1$: $f(x,1) = 1 \cdot e^{-1(5-x)} = e^{x-5}$ for $0 \le x \le 4$.
An increasing exponential curve starting at $e^{-5} \approx 0.0067$ mg/liter at $x=0$ and rising to $e^{-1} \approx 0.368$ mg/liter at $x=4$.

\textbf{Answer.}
(a) $f(4,t) = t e^{-t}$: tracks drug concentration over time following a $4$ mg dose; starts at $0$, peaks at $1/e \approx 0.368$ mg/liter at $t=1$ hour, and decays toward $0$.
(b) $f(x,1) = e^{x-5}$: shows drug concentration $1$ hour after injection as a function of dose; concentration increases exponentially with dose size.
\fig{.75}{answer-12-2-21.png}

\subsection{12.2, Exercise 27: Identify the Fixed Variable}
\textbf{Question.} Without a calculator or computer, for $z = x^2 + 2xy^2$, determine which of (I)--(II) in Figure 12.30 are cross-sections with $x$ fixed and which are cross-sections with $y$ fixed.
\fig{.65}{source-12-2-27.png}

\textbf{Relevant method.} PT-05: Multi-Curve Families of Vertical Cross-Sections.

\textbf{Necessary work.}
Fix $x = a$: $z = a^2 + 2ay^2$. These are parabolas in $y$ symmetric about $y=0$, opening upward if $a > 0$ and downward if $a < 0$, with vertex at $(y,z) = (0, a^2)$. This matches Family I.
Fix $y = b$: $z = x^2 + 2b^2 x = (x + b^2)^2 - b^4$. These are parabolas in $x$ all opening upward, all passing through the origin $(x,z)=(0,0)$, with shifted vertices at $(-b^2, -b^4)$. This matches Family II.

\textbf{Answer.}
\textbf{I} has $x$ fixed; \textbf{II} has $y$ fixed.

\subsection{12.2, Exercise 31: Match Cross-Section Families}
\textbf{Question.} Without a calculator or computer, match the functions with their cross-sections with $x$ fixed in Figure 12.31:
\[
\begin{array}{ll}
\text{(a) } z = 1/(1+x^2+y^2), & \text{(b) } z = 1 + x + y, \\
\text{(c) } z = e^{-x+y},       & \text{(d) } z = e^{x-y}, \\
\text{(e) } z = \sin(xy),       & \text{(f) } z = x^2.
\end{array}
\]
\fig{.65}{source-12-2-31.png}

\textbf{Relevant method.} PT-05: Multi-Curve Families of Vertical Cross-Sections.

\textbf{Necessary work.}
Fix $x = a$:
\begin{itemize}[nosep]
    \item (a) $z = \frac{1}{(1+a^2)+y^2}$: symmetric bell-shaped humps centered at $y=0 \implies$ Family II.
    \item (b) $z = (1+a) + y$: parallel straight lines of slope $1 \implies$ Family I.
    \item (c) $z = e^{-a} e^y$: increasing exponential curves $\implies$ Family III.
    \item (d) $z = e^a e^{-y}$: decreasing exponential curves $\implies$ Family VI.
    \item (e) $z = \sin(ay)$: sine waves whose frequency varies with $a \implies$ Family V.
    \item (f) $z = a^2$: horizontal straight lines independent of $y \implies$ Family IV.
\end{itemize}

\textbf{Answer.}
(a) $\to$ II; (b) $\to$ I; (c) $\to$ III; (d) $\to$ VI; (e) $\to$ V; (f) $\to$ IV.

\subsection{12.2, Exercise 41: Travelling Wave in Canal}
\textbf{Question.} A wave travels along a canal. Let $x$ be the distance along the canal, $t$ be the time, and $z$ be the height of the water above the equilibrium level. The graph of $z$ as a function of $x$ and $t$ is in Figure 12.33.
\begin{enumerate}[label=(\alph*),nosep]
    \item Draw the profile of the wave for $t = -1, 0, 1, 2$. (Put the $x$-axis to the right and the $z$-axis vertical.)
    \item Is the wave traveling in the direction of increasing or decreasing $x$?
    \item Sketch a surface representing a wave traveling in the opposite direction.
\end{enumerate}
\fig{.42}{source-12-2-41.png}

\textbf{Relevant method.} PT-05: Multi-Curve Families of Vertical Cross-Sections.

\textbf{Necessary work.}
(a) Fixed time $t$ gives snapshots of wave profile $z$ versus distance $x$.
(b) In Figure 12.33, as $t$ increases, the crest of the wave shifts toward larger values of $x$. Hence, the wave travels in the direction of increasing $x$.
(c) A wave in the opposite direction has ridges moving toward decreasing $x$ as $t$ increases (modeled by $z = g(x + vt)$).

\textbf{Answer.}
(a) Four hump-shaped curves translating toward larger $x$ as $t$ increases from $-1$ to $2$.
(b) Direction of \textbf{increasing $x$}.
(c) Surface with crests sloping in the opposite direction (sketched below).
\fig{.75}{answer-12-2-41.png}
\fig{.55}{answer-12-2-41c.png}

\subsection{12.3, Exercise 1: Surface to Contours (Downward Bowl)}
\textbf{Question.} Sketch a possible contour diagram for the surface below, marked with reasonable $z$-values. (Note: There are many possible answers.)
\fig{.32}{source-12-3-1.png}

\textbf{Relevant method.} PT-08: Algebraic Contours and Spacing.

\textbf{Necessary work.}
The surface is a downward-opening radially symmetric bowl with a maximum peak in the center. Horizontal slices are concentric circles. Since the peak is highest in the center, contour labels must decrease from the inside outward. A consistent mathematical model is $z = 4 - x^2 - y^2$.

\textbf{Answer.}
Concentric circles with labels decreasing outward (e.g., $c = 3, 2, 1, 0$ with central peak $4$).
\fig{.45}{answer-12-3-1.png}

\subsection{12.3, Exercise 3: Surface to Contours (Circular Depression)}
\textbf{Question.} Sketch a possible contour diagram for the surface below, marked with reasonable $z$-values. (Note: There are many possible answers.)
\fig{.32}{source-12-3-3.png}

\textbf{Relevant method.} PT-08: Algebraic Contours and Spacing.

\textbf{Necessary work.}
The surface is a circular depression (crater/basin) with a minimum at the center, rising outward and flattening asymptotically. Horizontal slices are concentric circles whose labels must increase from the center outward. A representative model is $z = 4(1 - e^{-(x^2+y^2)})$.

\textbf{Answer.}
Concentric circles with labels increasing outward from the center (e.g., center $0$, followed by $1, 2, 3, 3.5$).
\fig{.45}{answer-12-3-3.png}

\subsection{12.3, Exercise 5: Read a Contour Value}
\textbf{Question.} Use the contour diagram of $f(x,y)$ given in Figure 12.52 to find $f(1,3)$.
\fig{.45}{source-12-3-5.png}

\textbf{Relevant method.} PT-10: Read Labelled Contour Map.

\textbf{Necessary work.}
Locate coordinate $x = 1$ on the horizontal axis and $y = 3$ on the vertical axis. The point $(1,3)$ lies directly on the contour line labelled $-1$.

\textbf{Answer.}
$\boxed{f(1,3) = -1}$.

\subsection{12.3, Exercise 7: Find a Point on a Level Set}
\textbf{Question.} Use the contour diagram of $f(x,y)$ given in Figure 12.52 to find a point where $f(x,y) = 0$.

\textbf{Relevant method.} PT-10: Read Labelled Contour Map.

\textbf{Necessary work.}
Trace the contour line labelled $0$. The curve passes through the origin $(0,0)$. Other points along this zero contour include approximately $(2, 1.5)$ and $(-2, -1.5)$.

\textbf{Answer.}
$\boxed{(0,0)}$ (any coordinates on the $c=0$ curve are acceptable).

\subsection{12.3, Exercise 9: Contour Through a Specific Point}
\textbf{Question.} Let $f(x,y) = 3x^2 y + 7x + 20$. Find an equation for the contour that goes through the point $(5,10)$.

\textbf{Relevant method.} PT-08: Algebraic Contours.

\textbf{Necessary work.}
First compute the output value at the point $(5,10)$:
\[ c = f(5,10) = 3(5^2)(10) + 7(5) + 20 = 3(25)(10) + 35 + 20 = 750 + 55 = 805. \]
Set $f(x,y) = c$:
\[ 3x^2 y + 7x + 20 = 805 \implies 3x^2 y + 7x = 785. \]

\textbf{Answer.}
$\boxed{3x^2 y + 7x + 20 = 805}$ (or equivalently, $y = \frac{785 - 7x}{3x^2}$ for $x \ne 0$).

\subsection{12.3, Exercise 11: Match Surfaces with Contour Diagrams}
\textbf{Question.} Match the surfaces (a)--(e) in Figure 12.53 with the contour diagrams (I)--(V) in Figure 12.54.
\fig{.60}{source-12-3-11.png}

\textbf{Relevant method.} PT-08: Contours and Spacing; PT-06: Cylinders.

\textbf{Necessary work.}
\begin{itemize}[nosep]
    \item Surface (a) is a downward circular bowl $\implies$ Concentric circles with values decreasing outward $\implies$ Diagram III.
    \item Surface (b) is a trough parallel to the $x$-axis, minimum along the central line $\implies$ Paired horizontal lines with values increasing away from center $\implies$ Diagram I.
    \item Surface (c) is a downward parabolic cylinder parallel to the $y$-axis $\implies$ Paired vertical lines with values decreasing away from central ridge $\implies$ Diagram V.
    \item Surface (d) is an upward cone with constant slope $\implies$ Equally spaced concentric circles increasing outward $\implies$ Diagram II.
    \item Surface (e) is a high central circular mound $\implies$ Concentric circles with highest values in center $\implies$ Diagram IV.
\end{itemize}

\textbf{Answer.}
(a) $\to$ III; (b) $\to$ I; (c) $\to$ V; (d) $\to$ II; (e) $\to$ IV.

\subsection{12.3, Exercise 13: Tables to Contour Diagrams}
\textbf{Question.} Match Tables 12.6--12.9 with contour diagrams (I)--(IV) in Figure 12.56.
\fig{.55}{source-12-3-13.png}

\textbf{Relevant method.} PT-11: Data Tables and Contours.

\textbf{Necessary work.}
\begin{itemize}[nosep]
    \item Table 12.6: Values increase symmetrically away from a central minimum at $(0,0)$ $\implies$ Upward bowl $\implies$ Diagram II.
    \item Table 12.7: Values decrease symmetrically away from a central maximum at $(0,0)$ $\implies$ Downward bowl $\implies$ Diagram III.
    \item Table 12.8: Entries depend only on $x$ (columns constant, rows identical) $\implies$ Vertical contour lines $\implies$ Diagram IV.
    \item Table 12.9: Entries depend only on $y$ (rows constant, columns identical) $\implies$ Horizontal contour lines $\implies$ Diagram I.
\end{itemize}

\textbf{Answer.}
Table 12.6 $\to$ II; Table 12.7 $\to$ III; Table 12.8 $\to$ IV; Table 12.9 $\to$ I.

\subsection{12.3, Exercise 15: Parallel Line Contours}
\textbf{Question.} Sketch a contour diagram for $f(x,y) = 3x + 3y$ with at least four labeled contours. Describe in words the contours and how they are spaced.

\textbf{Relevant method.} PT-09: Contours of Linear Functions.

\textbf{Necessary work.}
Set $3x + 3y = c \implies y = -x + c/3$.
Select $c = -6, -3, 0, 3, 6$:
$y = -x - 2, \ -x - 1, \ -x, \ -x + 1, \ -x + 2$.
Perpendicular distance between adjacent lines for $\Delta c = 3$ is $D = \frac{|\Delta c|}{\sqrt{3^2+3^2}} = \frac{3}{3\sqrt{2}} = \frac{1}{\sqrt{2}}$.

\textbf{Answer.}
Parallel straight lines with slope $-1$, equally spaced for equal increments of $c$.
\fig{.45}{answer-12-3-15.png}

\subsection{12.3, Exercise 17: Downward Bowl Contours}
\textbf{Question.} Sketch a contour diagram for $f(x,y) = -x^2 - y^2 + 1$ with at least four labeled contours. Describe in words the contours and how they are spaced.

\textbf{Relevant method.} PT-08: Algebraic Contours and Spacing.

\textbf{Necessary work.}
Set $1 - x^2 - y^2 = c \implies x^2 + y^2 = 1 - c$.
Requires $c \le 1$. Choose $c = 1, 0, -1, -2, -3$:
$r = 0, \ 1, \ \sqrt{2} \approx 1.41, \ \sqrt{3} \approx 1.73, \ 2.0$.
Radial gaps: $1.0, 0.41, 0.32, 0.27$.

\textbf{Answer.}
Concentric circles centered at $(0,0)$ with labels decreasing outward. The distance between successive circles for equal $\Delta c$ decreases outward, showing increasing steepness.
\fig{.45}{answer-12-3-17.png}

\subsection{12.3, Exercise 19: Parabolic Contours}
\textbf{Question.} Sketch a contour diagram for $f(x,y) = y - x^2$ with at least four labeled contours. Describe in words the contours and how they are spaced.

\textbf{Relevant method.} PT-08: Algebraic Contours; PT-04: Parabolic Slices.

\textbf{Necessary work.}
Set $y - x^2 = c \implies y = x^2 + c$.
Choose $c = -2, -1, 0, 1, 2$.
Each curve is an upward-opening parabola with vertex $(0,c)$. Along any vertical line $x = x_0$, the vertical spacing $\Delta y = \Delta c$ is strictly constant.

\textbf{Answer.}
Upward-opening parabolas with vertices at $(0,c)$. Vertical spacing between adjacent curves is constant, but perpendicular separation shrinks as $|x|$ increases.
\fig{.45}{answer-12-3-19.png}

\subsection{12.3, Exercise 25: Interpret Economic Satisfaction Contours}
\textbf{Question.} Each contour diagram (a)--(c) in Figure 12.59 shows satisfaction with quantities of two items $X$ and $Y$ combined. Match (a)--(c) with the items in (I)--(III):
\begin{enumerate}[label=(\Roman*),nosep]
    \item $X$: Income; $Y$: Leisure time.
    \item $X$: Income; $Y$: Hours worked.
    \item $X$: Hours worked; $Y$: Time spent commuting.
\end{enumerate}
\fig{.55}{source-12-3-25.png}

\textbf{Relevant method.} PT-10: Read Contour Maps; PT-14: Economic Models.

\textbf{Necessary work.}
\begin{itemize}[nosep]
    \item In (a), satisfaction increases with $X$ (rightward) and decreases with $Y$ (upward). Matches II: income is desirable, hours worked is undesirable.
    \item In (b), satisfaction increases with both $X$ (rightward) and $Y$ (upward). Matches I: both income and leisure time are desirable.
    \item In (c), satisfaction decreases with both $X$ (rightward) and $Y$ (upward). Matches III: both hours worked and commute time are undesirable.
\end{itemize}

\textbf{Answer.}
(a) $\to$ II; (b) $\to$ I; (c) $\to$ III.

\subsection{12.3, Exercise 27: Trade-off Along Dan's Happiness Contour}
\textbf{Question.} Figure 12.61 shows a contour diagram of Dan's happiness with snacks of different numbers of cherries and grapes. (a) What is the slope of the contours? (b) What does the slope tell you?
\fig{.42}{source-12-3-27.png}

\textbf{Relevant method.} PT-14: Economic Trade-offs along Contours.

\textbf{Necessary work.}
(a) On any contour line, increasing cherries by $1$ unit corresponds to decreasing grapes by $2$ units:
\[ \text{Slope} = \frac{\Delta(\text{grapes})}{\Delta(\text{cherries})} = \frac{-2}{1} = -2. \]
(b) Moving along a contour keeps happiness constant. A slope of $-2$ means $1$ cherry compensates for the loss of $2$ grapes.

\textbf{Answer.}
(a) Slope is $-2$ grapes per cherry.
(b) Dan is equally happy if he trades $2$ grapes for $1$ cherry; his marginal rate of substitution is $2$ grapes per cherry.

\subsection{12.3, Exercise 47: Manufacturer Profit Contours}
\textbf{Question.} A manufacturer sells two goods, one at a price of \$3000 a unit and the other at a price of \$12,000 a unit. A quantity $q_1$ of the first good and $q_2$ of the second good are sold at a total cost of \$4000 to the manufacturer.
\begin{enumerate}[label=(\alph*),nosep]
    \item Express the manufacturer's profit, $\pi$, as a function of $q_1$ and $q_2$.
    \item Sketch curves of constant profit in the $q_1q_2$-plane for $\pi = 10{,}000$, $\pi = 20{,}000$, and $\pi = 30{,}000$, and the break-even curve $\pi = 0$.
\end{enumerate}

\textbf{Relevant method.} PT-09: Linear Contours; PT-14: Economic Models.

\textbf{Necessary work.}
(a) Profit = Revenue $-$ Cost:
\[ \pi(q_1, q_2) = 3000 q_1 + 12000 q_2 - 4000 \text{ (dollars)}. \]
(b) Set $\pi = c$:
\[ 3000 q_1 + 12000 q_2 = c + 4000 \implies q_2 = -\frac{1}{4}q_1 + \frac{c + 4000}{12000}. \]
Since $q_1, q_2 \ge 0$, contours are parallel lines of slope $-1/4$ restricted to the first quadrant.
\begin{itemize}[nosep]
    \item $\pi = 0$: $q_1$-intercept is $4/3 \approx 1.33$, $q_2$-intercept is $1/3 \approx 0.33$.
    \item $\pi = 10{,}000$: $q_1$-intercept is $14/3 \approx 4.67$, $q_2$-intercept is $7/6 \approx 1.17$.
    \item $\pi = 20{,}000$: $q_1$-intercept is $8$, $q_2$-intercept is $2$.
    \item $\pi = 30{,}000$: $q_1$-intercept is $34/3 \approx 11.33$, $q_2$-intercept is $17/6 \approx 2.83$.
\end{itemize}

\textbf{Answer.}
(a) $\boxed{\pi(q_1, q_2) = 3000q_1 + 12000q_2 - 4000}$ dollars.
(b) Parallel line segments in the first quadrant of slope $-1/4$, shifting outward as profit increases.
\fig{.65}{answer-12-3-47.png}

\clearpage

% =========================================================================
% SECTION D: ORIGINAL SELF-TEST QUESTIONS
% =========================================================================
\section{Original Self-Test Questions}\label{sec:selftest-questions}
Attempt all 15 questions independently before checking answers in Section~\ref{sec:selftest-answers}.

\begin{enumerate}[label=\textbf{Q\arabic*.},leftmargin=*]
    \item A set in $\R^3$ is defined by $x^2 + y^2 + z^2 = 9$. Is the entire set the graph of a single function $z = f(x,y)$? Justify using an explicit input point, and state the exact natural domain after splitting the surface into two function branches.
    
    \item Describe the geometric surface $z = 3 - 2x + y$ in $\R^3$: name the object, determine all coordinate axis intercepts, state the $x=0$ and $y=0$ vertical traces, specify directions of increase, and state whether it has finite or infinite extent.
    
    \item For the hyperbolic paraboloid $z = x^2 - y^2$, determine the equations of the cross-sections in the planes $x = 2$ and $y = -1$. Identify the geometric shape, opening direction, and containing plane for each curve.
    
    \item For the function $F(x,y) = y^2 + 2xy$, write the algebraic equations of the vertical cross-sections for $x = 0$, $x = 1$, $y = 0$, and $y = 2$. State the designated horizontal plotting coordinate for each family.
    
    \item Describe the surface defined by $y^2 + z^2 = 4$ in $\R^3$. Which coordinate is free? State the axis of extrusion, and evaluate whether the surface is the graph of a single-valued function $z = f(x,y)$.
    
    \item For $f(x,y) = x^2 + 4y^2$, determine the permissible values of $c$, find the contour equations for $c = 0, 4, 16$, and compute their axis intercepts. What is the 3D shape of the surface?
    
    \item For the downward cone $g(x,y) = 9 - \sqrt{x^2+y^2}$, determine the radii of the contours corresponding to $g = 8, 7, 6$. Explain whether function values increase toward or away from the origin.
    
    \item For the linear function $p(x,y) = 3x - y + 2$, write the contour line equations for $p = -1, 2, 5$ in slope-intercept form. State their common slope and a vector direction in the $xy$-plane along which $p$ increases.
    
    \item The table below records values of $v(x,y)$ in metres on $[0,2] \times [0,2]$.
    \begin{center}
    \begin{tabular}{c|rrr}
    \toprule
    $y \backslash x$ (m) & $0$ & $1$ & $2$ \\
    \midrule
    $0$ & $9$ & $7$ & $5$ \\
    $1$ & $8$ & $6$ & $4$ \\
    $2$ & $7$ & $5$ & $3$ \\
    \bottomrule
    \end{tabular}
    \end{center}
    Find $v(2,1)$ exactly. Using linear interpolation along the rows, estimate where the contour $v = 5$ crosses $y = 0$, $y = 1$, and $y = 2$. Fit a linear equation to these points.
    
    \item On the contour map of $q(x,y) = 12 - x^2 - y^2$ below, point $A = (1.5, 0)$ lies between the contours $q = 10$ and $q = 8$. Give a strict bounding bracket for $q(A)$, estimate its value via linear interpolation, and compute its exact value using the formula. What do closely spaced contours signify on this map?
    \fig{.65}{10-self-test.png}
    
    \item For the saddle function $s(x,y) = x^2 - y^2$, describe the geometric structure of the level sets $s = 0$, $s = 3$, and $s = -3$. How many disconnected branches does each level set possess? Explain why the intersection of contours at $(0,0)$ does not violate the non-intersecting rule.
    
    \item A triangular camping tent covers the ground rectangle $0 \le x \le 4$ m and $0 \le y \le 2$ m. Its horizontal ridge lies along $x = 2$ at height $6$ m, and both side ground edges ($x=0$ and $x=4$) have height $0$. Formulate the piecewise roof height $h(x,y)$ and determine the equation of the $3$ m contour.
    
    \item \textbf{(Supplementary)} Starting from the standard circular paraboloid $z = x^2+y^2$, describe the geometric transformations that produce $z = (x+2)^2 + (y-1)^2 - 3$. Identify its vertex coordinates, opening direction, and range.
    
    \item \textbf{(Supplementary)} A production process is modeled by $P(N,V) = \sqrt{NV}$ for labor $N > 0$ and capital $V > 0$. Derive the algebraic formula for the isoquant $P = 4$. If labor is doubled, by what factor must capital change to keep production constant?
    
    \item \textbf{(Preview)} For the three-variable function $H(x,y,z) = x^2 + y^2 + z^2$, describe the level surfaces corresponding to $H = 0, 1, 4$. For $J(x,y,z) = z - y$, describe the level surface $J = 2$. Identify the geometric classification of each surface.
\end{enumerate}

\clearpage

% =========================================================================
% SECTION E: SELF-TEST ANSWERS
% =========================================================================
\section{Self-Test Answers}\label{sec:selftest-answers}
Each answer references its corresponding Problem Type ID from \texttt{Study-Notes.tex}.

\begin{enumerate}[label=\textbf{A\arabic*.},leftmargin=*]
    \item \textbf{Reference: PT-02.} No, the entire sphere fails the vertical-line test because the line $x=0, y=0$ yields $z = \pm 3$, intersecting at $(0,0,3)$ and $(0,0,-3)$. Splitting into two single-valued functions yields $z = +\sqrt{9-x^2-y^2}$ and $z = -\sqrt{9-x^2-y^2}$, each with natural domain $D = \{(x,y) \in \R^2 : x^2+y^2 \le 9\}$.
    
    \item \textbf{Reference: PT-01.} Infinite plane. Intercepts are $(3/2, 0, 0)$, $(0, -3, 0)$, and $(0, 0, 3)$. Traces are $z = 3 + y$ in the plane $x=0$, and $z = 3 - 2x$ in the plane $y=0$. Height increases with decreasing $x$ and increasing $y$. The surface has infinite, unbounded extent.
    
    \item \textbf{Reference: PT-04.} For $x = 2$: $z = 4 - y^2$ in the vertical plane $x = 2$, a downward-opening parabola with vertex $(2, 0, 4)$. For $y = -1$: $z = x^2 - 1$ in the vertical plane $y = -1$, an upward-opening parabola with vertex $(0, -1, -1)$.
    
    \item \textbf{Reference: PT-05.} Cross-sections with $x$ fixed (plotted on $yz$-axes): $x=0 \implies z = y^2$; $x=1 \implies z = y^2 + 2y$. Cross-sections with $y$ fixed (plotted on $xz$-axes): $y=0 \implies z = 0$; $y=2 \implies z = 4 + 4x$.
    
    \item \textbf{Reference: PT-06.} Variable $x$ is free. The surface is a right circular cylinder of radius $2$ centered on the $x$-axis. It fails the vertical-line test because vertical lines through interior points of $|y|<2$ intersect the cylinder twice ($z = \pm\sqrt{4-y^2}$); hence, it is not a single function $z=f(x,y)$.
    
    \item \textbf{Reference: PT-08.} Requires $c \ge 0$. For $c=0$: single point $(0,0)$. For $c=4$: ellipse $x^2 + 4y^2 = 4 \implies x^2/4 + y^2 = 1$, intercepts $(\pm 2, 0)$ and $(0, \pm 1)$. For $c=16$: ellipse $x^2/16 + y^2/4 = 1$, intercepts $(\pm 4, 0)$ and $(0, \pm 2)$. The 3D surface is an elliptic paraboloid opening upward.
    
    \item \textbf{Reference: PT-08.} $9 - r = c \implies r = 9 - c$. Radii for $g = 8, 7, 6$ are $r = 1, 2, 3$, respectively. Function values increase toward the origin, where the peak value $g(0,0)=9$ is attained.
    
    \item \textbf{Reference: PT-09.} Setting $3x - y + 2 = c \implies y = 3x + (2-c)$.
    For $p = -1$: $y = 3x + 3$.
    For $p = 2$: $y = 3x$.
    For $p = 5$: $y = 3x - 3$.
    Common slope is $3$. A vector direction of increasing $p$ is $(3, -1)$ (or any vector where $3\Delta x - \Delta y > 0$).
    
    \item \textbf{Reference: PT-11.} Exact lookup: row $y=1$, column $x=2$ gives $v(2,1) = 4$ m.
    Interpolation for $v=5$: at $y=0$, $x=2$ exactly; at $y=1$, $x = 1 + \frac{5-6}{4-6}(2-1) = 1.5$; at $y=2$, $x=1$ exactly.
    Points are $(2,0)$, $(1.5,1)$, and $(1,2)$. Connecting them yields the line $2x + y = 4$ on $[0,2]^2$.
    
    \item \textbf{Reference: PT-10.} Bounding bracket: $8 < q(A) < 10$. Linear radial interpolation yields $q_{\text{est}} \approx 9.71$. Exact substitution: $q(1.5,0) = 12 - (1.5)^2 = 9.75$. Closely spaced contours indicate rapid elevation change over short horizontal distance (steep surface).
    
    \item \textbf{Reference: PT-12.} $s = 0$ consists of two intersecting lines $y = \pm x$ (one connected set). $s = 3$ is a hyperbola opening horizontally along the $x$-axis (two disconnected branches). $s = -3$ is a hyperbola opening vertically along the $y$-axis (two disconnected branches). Crossing at $(0,0)$ involves curves of the same height ($c=0$), which is fully consistent with the single-valued function definition.
    
    \item \textbf{Reference: PT-13.} Piecewise height: $h(x,y) = 3x$ on $0 \le x \le 2$, and $h(x,y) = 12 - 3x$ on $2 \le x \le 4$, with $0 \le y \le 2$. At $h = 3$ m: $3x = 3 \implies x = 1$ m, and $12 - 3x = 3 \implies x = 3$ m. The contour consists of two vertical line segments $x = 1$ and $x = 3$ spanning $0 \le y \le 2$.
    
    \item \textbf{Reference: PT-03.} Horizontal translation by $(-2, 1, 0)$ followed by vertical translation by $-3$. The surface is an upward circular paraboloid with vertex $(-2, 1, -3)$, axis $x=-2, y=1$, and range $[-3, \infty)$.
    
    \item \textbf{Reference: PT-14.} $\sqrt{NV} = 4 \implies NV = 16 \implies V = 16/N$ ($N>0$). If labor doubles ($N \mapsto 2N$), capital must be halved ($V \mapsto V/2$) to keep production at $4$.
    
    \item \textbf{Reference: PT-15.} For $H(x,y,z)$: $H=0$ is the single point $(0,0,0)$; $H=1$ is a sphere of radius $1$ centered at $(0,0,0)$; $H=4$ is a sphere of radius $2$ centered at $(0,0,0)$. For $J(x,y,z)$: $J=2$ is the infinite plane $z = y + 2$ parallel to the $x$-axis.
\end{enumerate}

\clearpage

% =========================================================================
% SECTION F: COMPACT COVERAGE AND UNRESOLVED-ITEMS APPENDIX
% =========================================================================
\section{Compact Coverage and Unresolved-Items Appendix}\label{sec:coverage-audit}

\subsection{Scope and Provenance}
The source materials comprise 54 pages across six weekly PDF documents in \texttt{Lec Notes} and Sections 12.2 and 12.3 of \emph{Calculus, Eighth Edition} (printed pp.~702--725). All 26 assigned textbook exercises (13 from Section 12.2 and 13 from Section 12.3) and all weekly lecture problems are solved in full. Supporting graphs were generated via \texttt{graphs/generate\_graphs.py} or cropped directly from source textbook figures.

\subsection{Compact Source Coverage Matrix}
\small
\begin{longtable}{p{.38\linewidth}p{.38\linewidth}p{.18\linewidth}}
\toprule
Source Document \& Pages & Mathematical Topic Covered & Location / PT ID \\
\midrule
\endhead
\texttt{Week2.pdf}, p.~1 & Plane surface and linear cross-sections & PT-01, PT-04, \S B.1 \\
\texttt{Week2.pdf}, p.~2 & Quadric surface catalog and cylinder test & PT-01, PT-06, \S B.2 \\
\texttt{Week2.pdf}, p.~3 & Multi-curve cross-sections ($y^3+xy$) & PT-05, \S B.3 \\
\texttt{Week2.pdf}, pp.~4--5 & Paraboloid and cone contours, spacing & PT-08, \S B.4 \\
\texttt{Week2.pdf}, p.~6 & Linear contours ($2y-x$) and gradient & PT-09, \S B.5 \\
\texttt{Week2.pdf}, pp.~7--8 & Saddle table ($x^2-y^2$) and hyperbolic levels & PT-12, \S B.6 \\
\texttt{Week2.pdf}, pp.~9--10 & Tent roof, bounded contours, 3D distance & PT-13, \S B.7 \\
\texttt{sep+14.pdf}, pp.~1--4 & Vertical-line test, shifts, saddle traces & PT-02, PT-03, \S B.8 \\
\texttt{sep+16.pdf}, pp.~1--4 & Free variables, cylinders, bowl spacing & PT-06, PT-08, \S B.9 \\
\texttt{sep+18.pdf}, pp.~1--4 & Table interpolation and surface contrasts & PT-11, PT-12, \S B.10 \\
\texttt{MAT235H-Sept-16}, pp.~2--4 & Functions of 2 variables, cylinder volume & PT-01, PT-06, \S B.11 \\
\texttt{MAT235H-Sept-16}, p.~7 & Atmospheric wind-chill table lookup & PT-11, \S B.12 \\
\texttt{MAT235H-Sept-16}, pp.~5,6,9,11 & 3D points, planes, heater cross-sections & PT-01, PT-05, \S B.13 \\
\texttt{MAT235H-Sept-16}, pp.~13--19 & Cylinders, sphere, hemisphere slice & PT-02, PT-04, \S B.14 \\
\texttt{MAT235H-Sept-16}, pp.~24--27 & Airline table, plane intercepts, linearity & PT-09, PT-11, \S B.15 \\
\texttt{MAT235H-Sept-17}, pp.~1--3 & Plane equations and intercept graphing & PT-01, PT-09, \S B.15 \\
\texttt{MAT235H-Sept-17}, pp.~4--5 & Level surfaces in $\R^3$ (ice block, $T=c$) & PT-15, \S B.16 \\
Textbook 12.2, pp.~702--707 & Theory: graphs, vertical-line test, traces & PT-01 to PT-06 \\
Textbook 12.2, pp.~707--711 & Exercises 1, 3, 5, 7, 9, 11, 13, 15, 17, 21, 27, 31, 41 & \S C (12.2 Exercises) \\
Textbook 12.3, pp.~711--718 & Theory: contours, steepness, Cobb-Douglas & PT-07 to PT-14 \\
Textbook 12.3, pp.~718--725 & Exercises 1, 3, 5, 7, 9, 11, 13, 15, 17, 19, 25, 27, 47 & \S C (12.3 Exercises) \\
\bottomrule
\end{longtable}
\normalsize

\subsection{Mathematical Corrections, Ambiguities, and Modeling Boundaries}
\begin{itemize}[leftmargin=*]
    \item \textbf{Slice Terminology Correction (\texttt{sep+14.pdf}, p.~4):} The handwritten lecture notes refer to a vertical slice $y = a$ of a circular paraboloid as a ``paraboloid.'' A vertical slice of a paraboloid is a two-dimensional \emph{parabola} contained in the plane $y=a$, not a paraboloid. The plane condition and correct 2D conic classification are maintained throughout these notes.
    \item \textbf{Hemisphere Domain Restriction (\texttt{MAT235H-Sept-16.pdf}, p.~17):} The slide notes state $-1 \le x \le 1, -1 \le y \le 1$ alongside the upper hemisphere formula $z = \sqrt{1-x^2-y^2}$. A rectangular domain contains points like $(1,1)$ where $1-x^2-y^2 = -1 < 0$. The exact domain is strictly the circular disk $x^2 + y^2 \le 1$.
    \item \textbf{Table Interpolation vs Global Identity (\texttt{sep+18.pdf}, p.~2):} While the discrete data points in the table are consistent with the affine function $f(x,y) = 12 - 2x - y$, finite grid samples establish an empirical linear interpolation model, not a proven global planar identity. Solutions explicitly state the linear interpolation assumption.
    \item \textbf{Missing Slide Prompt (\texttt{MAT235H-Sept-17.pdf}, p.~1):} The opening slide displays the equations $z = 2 - 2x + y$ and $z = 2 - x - 2y$ without an explicit verbal instruction. These have been interpreted as finding coordinate intercepts and describing the planes, without inventing unstated requirements.
    \item \textbf{Empirical Readings and Estimates:} In Section 12.3 Exercise 21 (corn-yield diagram) and lecture heater sketches, values not lying on a labelled line are presented as justified brackets and estimates rather than exact measurements.
\end{itemize}

\end{document}
